\documentclass[12pt]{article}

% === Кодировки и язык ===
\usepackage{cmap}
\usepackage[T2A]{fontenc}
\usepackage{mathtext}
\usepackage[utf8]{inputenc}
\usepackage[english, russian]{babel}

% === Цвета ===
\usepackage[dvipsnames]{xcolor}

% === Математика ===
\usepackage{amsmath, amssymb, wasysym, amsthm}
\usepackage{stmaryrd}
\usepackage{proof}

% === Графика и float'ы ===
\usepackage{wrapfig, float}
\usepackage{graphicx}

% === TikZ / PGFPlots ===
\usepackage{pgfplots}
\pgfplotsset{compat=newest}
\usepackage{pgfplotstable}
\usepackage{tikz, tikz-3dplot}
\usetikzlibrary{calc, patterns, angles, quotes, intersections, automata, positioning, arrows.meta}

% === Таблицы ===
\usepackage{tabularx}
\usepackage{xltabular}
\renewcommand\tabularxcolumn[1]{m{#1}}

% === Разметка ===
\usepackage{geometry}
\geometry{margin=0.5in}

% === Списки и прочее ===
\usepackage[shortlabels]{enumitem}
\usepackage{verbatim}

% ==========================================
% === ТЕОРЕМНЫЕ ОКРУЖЕНИЯ (amsthm) ===
% ==========================================
\theoremstyle{plain}
\newtheorem{theorem}{Теорема}[section]
\newtheorem{corollary}{Следствие}[section]
\newtheorem{lemma}{Лемма}[section]
\newtheorem{statement}{Утверждение}[section]
\newtheorem{axiom}{Аксиома}[section]
\newtheorem{law}{Закон}[section]

\theoremstyle{definition}
\newtheorem{definition}{Определение}[section]
\newtheorem{example}{Пример}[section]
\newtheorem{problem}{Задача}
\newtheorem{method}{Метод}[section]
\newtheorem{event}{Событие}[section]

\theoremstyle{remark}
\newtheorem{note}{Заметка}[section]
\newtheorem{property}{Свойство}[section]
\newtheorem{properties}{Свойства}[section]
\newtheorem{principle}{Правило}[section]
\newtheorem{person}{Личность}[section]
\newtheorem*{solution}{Решение}

% ==========================================
% === КРАСИВЫЕ РАМКИ (tcolorbox) ===
% ==========================================
\usepackage{tcolorbox}
\tcbuselibrary{skins, breakable, theorems}

% Базовые стили для разных типов блоков
\tcbset{
    % Стиль для теорем, лемм (Синий - строгий)
    thmstyle/.style={
        enhanced, breakable,
        frame hidden, boxrule=0pt,
        colback=blue!5!white, % Фон: 5% синего, 95% белого
        borderline west={2.5pt}{0pt}{blue!70!black}, % Вертикальная линия слева
        before skip=10pt, after skip=10pt,
        left=4pt, right=4pt, top=4pt, bottom=4pt
    },
    % Стиль для доказательств (фиолетовый - процесс рассуждения)
    proofstyle/.style={
        enhanced, breakable,
        frame hidden, boxrule=0pt,
        colback=Plum!5!white,
        borderline west={2.5pt}{0pt}{Plum!80!black},
        before skip=10pt, after skip=10pt,
        left=4pt, right=4pt, top=4pt, bottom=4pt,
        sharp corners
    },
    % Стиль для определений и аксиом (Зеленый - фундаментальный)
    defstyle/.style={
        enhanced, breakable,
        frame hidden, boxrule=0pt,
        colback=ForestGreen!5!white,
        borderline west={2.5pt}{0pt}{ForestGreen!80!black},
        before skip=10pt, after skip=10pt,
        left=4pt, right=4pt, top=4pt, bottom=4pt
    },
    % Стиль для примеров и задач (Оранжевый - практический)
    exstyle/.style={
        enhanced, breakable,
        frame hidden, boxrule=0pt,
        colback=orange!5!white,
        borderline west={2.5pt}{0pt}{orange!90!black},
        before skip=10pt, after skip=10pt,
        left=4pt, right=4pt, top=4pt, bottom=4pt
    },
    % Стиль для заметок и свойств (Серый - нейтральный)
    remstyle/.style={
        enhanced, breakable,
        frame hidden, boxrule=0pt,
        colback=gray!5!white,
        borderline west={2.5pt}{0pt}{gray!70!black},
        before skip=10pt, after skip=10pt,
        left=4pt, right=4pt, top=4pt, bottom=4pt
    }
}

% Применяем стили к существующим окружениям amsthm
% 1. Строгая математика
\tcolorboxenvironment{theorem}{thmstyle}
\tcolorboxenvironment{lemma}{thmstyle}
\tcolorboxenvironment{corollary}{thmstyle}
\tcolorboxenvironment{statement}{thmstyle}
\tcolorboxenvironment{law}{thmstyle}

% 2. Доказательство
\tcolorboxenvironment{proof}{proofstyle}
\renewcommand{\qed}{}

% 3. Определения
\tcolorboxenvironment{definition}{defstyle}
\tcolorboxenvironment{axiom}{defstyle}

% 4. Практика
\tcolorboxenvironment{example}{exstyle}
\tcolorboxenvironment{problem}{exstyle}
\tcolorboxenvironment{method}{exstyle}
\tcolorboxenvironment{event}{exstyle}

% 5. Заметки и прочее
\tcolorboxenvironment{note}{remstyle}
\tcolorboxenvironment{property}{remstyle}
\tcolorboxenvironment{properties}{remstyle}
\tcolorboxenvironment{principle}{remstyle}
\tcolorboxenvironment{person}{remstyle}
\tcolorboxenvironment{solution}{remstyle}

% === Гиперссылки (почти последним!) ===
\usepackage[
colorlinks=true,
linkcolor=blue,
urlcolor=blue,
citecolor=blue,
bookmarks=true,
bookmarksopen=false
]{hyperref}


% === Дополнительные пакеты ===
\usepackage{multicol}
\usepackage{array}
\usepackage{booktabs}
\usepackage{tcolorbox}
\tcbuselibrary{theorems, skins, breakable}
\usetikzlibrary{graphs, positioning, arrows.meta, decorations.pathmorphing}

\begin{document}
\tableofcontents
\setcounter{tocdepth}{3}
\newpage

%%%%%%%%%%%%%%%%%%%%%%%%%%%%%%%%%%%%%%%%%%%%%%%%%%%%%%%%%%%%%
\part{Логика и Теория множеств}
%%%%%%%%%%%%%%%%%%%%%%%%%%%%%%%%%%%%%%%%%%%%%%%%%%%%%%%%%%%%%

\section{Высказывания. Основные логические операции}

\begin{definition}
\textbf{Высказывание} --- это утверждение, которому можно однозначно
приписать одно из двух значений истинности: \emph{истина} (1, И, T)
или \emph{ложь} (0, Л, F).
\end{definition}

\begin{example}
«$2 + 2 = 4$» --- истинное высказывание.
«$2 + 2 = 5$» --- ложное высказывание.
«Закройте дверь!» --- \emph{не} высказывание (вопросы и просьбы).
\end{example}

\begin{definition}
\textbf{Парадокс} --- высказывание, которое не является ни истинным,
ни ложным, или одновременно является и тем и другим.
\end{definition}

\begin{example}
\textbf{Парадокс лжеца:} «Данное высказывание ложно». Если оно истинно,
то оно ложно; если ложно, то истинно --- противоречие.

\textbf{Парадокс Рассела:} Рассмотрим множество $R$ всех множеств,
не содержащих себя в качестве элемента:
$R = \{x \mid x \notin x\}$.
Тогда $R \in R \Leftrightarrow R \notin R$ --- противоречие.
\end{example}

\begin{definition}
\textbf{Тавтология} --- формула, принимающая значение \emph{истина}
при любом наборе значений переменных.
\end{definition}

\begin{example}
$A \vee \neg A$ --- тавтология (закон исключённого третьего).
\end{example}

\begin{definition}
\textbf{Противоречие} (контрадикция) --- формула, принимающая значение
\emph{ложь} при любом наборе значений переменных.
\end{definition}

\begin{example}
$A \wedge \neg A$ --- противоречие.
\end{example}

\subsection{Основные логические операции}

\begin{definition}
\textbf{Отрицание} (инверсия) $\neg A$:
\[
\neg A = \begin{cases} 1, & A = 0 \\ 0, & A = 1 \end{cases}
\]
\end{definition}

\begin{definition}
\textbf{Конъюнкция} (логическое И) $A \wedge B$ (также $A \,\&\, B$,
$A \cdot B$):
\[
A \wedge B = 1 \iff A = 1 \text{ и } B = 1.
\]
\end{definition}

\begin{definition}
\textbf{Дизъюнкция} (логическое ИЛИ) $A \vee B$:
\[
A \vee B = 0 \iff A = 0 \text{ и } B = 0.
\]
\end{definition}

%-----------------------------------------
\section{Таблица истинности и логические схемы для формул}
%-----------------------------------------

\begin{definition}
\textbf{Таблица истинности} --- таблица, строки которой соответствуют
всем возможным наборам значений переменных (для $n$ переменных
таблица имеет $2^n$ строк), а столбцы --- значениям формул.
\end{definition}

Таблица истинности для основных операций:

\begin{center}
\begin{tabular}{|c|c||c|c|c|c|c|}
\hline
$A$ & $B$ & $\neg A$ & $A \wedge B$ & $A \vee B$ &
$A \to B$ & $A \leftrightarrow B$ \\
\hline
0 & 0 & 1 & 0 & 0 & 1 & 1 \\
0 & 1 & 1 & 0 & 1 & 1 & 0 \\
1 & 0 & 0 & 0 & 1 & 0 & 0 \\
1 & 1 & 0 & 1 & 1 & 1 & 1 \\
\hline
\end{tabular}
\end{center}

\begin{definition}
\textbf{Логическая схема} --- представление формулы в виде графа,
где узлы --- логические элементы (вентили): NOT, AND, OR,
а рёбра --- провода (передача значений).
\end{definition}

\begin{example}
Формула $(A \wedge B) \vee \neg C$ реализуется схемой:
входы $A, B$ подаются на вентиль AND; вход $C$ --- на NOT;
выходы AND и NOT подаются на вентиль OR.
\end{example}

%-----------------------------------------
\section{Свойства логических операций}\label{sec:logic-props}
%-----------------------------------------

\begin{property}[Идемпотентность]
\[
A \wedge A = A, \qquad A \vee A = A.
\]
\end{property}

\begin{property}[Коммутативность]
\[
A \wedge B = B \wedge A, \qquad A \vee B = B \vee A.
\]
\end{property}

\begin{property}[Ассоциативность]
\[
(A \wedge B) \wedge C = A \wedge (B \wedge C),
\qquad
(A \vee B) \vee C = A \vee (B \vee C).
\]
\end{property}

\begin{property}[Дистрибутивность]
\[
A \wedge (B \vee C) = (A \wedge B) \vee (A \wedge C),
\]
\[
A \vee (B \wedge C) = (A \vee B) \wedge (A \vee C).
\]
\end{property}

\begin{property}[Законы де Моргана]
\[
\neg(A \wedge B) = \neg A \vee \neg B,
\qquad
\neg(A \vee B) = \neg A \wedge \neg B.
\]
\end{property}

\begin{property}[Законы поглощения]
\[
A \wedge (A \vee B) = A,
\qquad
A \vee (A \wedge B) = A.
\]
\end{property}

\begin{property}[Законы нейтральных элементов]
\[
A \wedge 1 = A, \quad A \wedge 0 = 0,
\qquad
A \vee 0 = A, \quad A \vee 1 = 1.
\]
\end{property}

\begin{property}[Двойное отрицание]
\[
\neg \neg A = A.
\]
\end{property}

\begin{property}[Закон исключённого третьего]
\[
A \vee \neg A = 1.
\]
\end{property}

\begin{property}[Закон противоречия]
\[
A \wedge \neg A = 0.
\]
\end{property}

%-----------------------------------------
\section{Дистрибутивность, ассоциативность, коммутативность}
%-----------------------------------------

Эти свойства уже перечислены в \S\ref{sec:logic-props}. Приведём \emph{доказательства}
перебором таблицы истинности для наиболее важных.

\begin{theorem}[Дистрибутивность $\wedge$ над $\vee$]
$A \wedge (B \vee C) = (A \wedge B) \vee (A \wedge C)$.
\end{theorem}
\begin{proof}
Проверяем все 8 наборов:
\begin{center}
\begin{tabular}{|c|c|c|c|c|c|c|}
\hline
$A$ & $B$ & $C$ & $B \vee C$ & $A \wedge(B\vee C)$ &
$(A\wedge B)\vee(A\wedge C)$ \\
\hline
0&0&0& 0 & 0 & 0 \\
0&0&1& 1 & 0 & 0 \\
0&1&0& 1 & 0 & 0 \\
0&1&1& 1 & 0 & 0 \\
1&0&0& 0 & 0 & 0 \\
1&0&1& 1 & 1 & 1 \\
1&1&0& 1 & 1 & 1 \\
1&1&1& 1 & 1 & 1 \\
\hline
\end{tabular}
\end{center}
Столбцы 5 и 6 совпадают.
\end{proof}

%-----------------------------------------
\section{Законы де Моргана от нескольких аргументов}
%-----------------------------------------

\begin{theorem}[Обобщённый закон де Моргана]
Для любого натурального $n \geq 2$:
\[
\neg(A_1 \wedge A_2 \wedge \cdots \wedge A_n)
= \neg A_1 \vee \neg A_2 \vee \cdots \vee \neg A_n,
\]
\[
\neg(A_1 \vee A_2 \vee \cdots \vee A_n)
= \neg A_1 \wedge \neg A_2 \wedge \cdots \wedge \neg A_n.
\]
\end{theorem}
\begin{proof}
Индукция по $n$.

\textbf{База} $n = 2$: классический закон де Моргана (проверяется таблицей).

\textbf{Шаг.} Пусть утверждение верно для $n-1$. Тогда:
\begin{align*}
\neg(A_1 \wedge \cdots \wedge A_n)
&= \neg\bigl((A_1 \wedge \cdots \wedge A_{n-1}) \wedge A_n\bigr) \\
&= \neg(A_1 \wedge \cdots \wedge A_{n-1}) \vee \neg A_n
   & &\text{(де Морган для 2 аргументов)}\\
&= (\neg A_1 \vee \cdots \vee \neg A_{n-1}) \vee \neg A_n
   & &\text{(предположение индукции)}\\
&= \neg A_1 \vee \cdots \vee \neg A_n.
\end{align*}
Аналогично для дизъюнкции.
\end{proof}

%-----------------------------------------
\section{Импликация и эквиваленция и их свойства}
%-----------------------------------------

\begin{definition}
\textbf{Импликация} $A \to B$ (\emph{«если $A$, то $B$»}):
\[
A \to B \stackrel{\mathrm{def}}{=} \neg A \vee B.
\]
Она ложна тогда и только тогда, когда $A = 1$ и $B = 0$.
\end{definition}

\begin{property}[Свойства импликации]
\begin{enumerate}[(i)]
\item \textbf{Контрапозиция:} $A \to B \equiv \neg B \to \neg A$.
\item \textbf{Транзитивность:} $(A \to B) \wedge (B \to C) \to (A \to C)$ ---
      тавтология.
\item \textbf{Рефлексивность:} $A \to A$ --- тавтология.
\item \textbf{Разложение:} $A \to B \equiv \neg A \vee B$.
\item $(A \to B) \equiv (\neg B \to \neg A)$ (контрапозитив).
\item $A \to (B \to C) \equiv (A \wedge B) \to C$ (экспортация).
\end{enumerate}
\end{property}

\begin{proof}[Доказательство контрапозиции]
\[
\neg B \to \neg A
= \neg(\neg B) \vee \neg A
= B \vee \neg A
= \neg A \vee B
= A \to B. \qquad \square
\]
\end{proof}

\begin{definition}
\textbf{Эквиваленция} (биконъюнкция) $A \leftrightarrow B$
(\emph{«$A$ тогда и только тогда, когда $B$»}):
\[
A \leftrightarrow B \stackrel{\mathrm{def}}{=}
(A \to B) \wedge (B \to A)
= (\neg A \vee B) \wedge (\neg B \vee A).
\]
Она истинна тогда и только тогда, когда $A$ и $B$ имеют одинаковые
значения.
\end{definition}

\begin{property}[Свойства эквиваленции]
\begin{enumerate}[(i)]
\item \textbf{Рефлексивность:} $A \leftrightarrow A$ --- тавтология.
\item \textbf{Симметричность:} $A \leftrightarrow B \equiv B \leftrightarrow A$.
\item \textbf{Транзитивность:}
      $(A \leftrightarrow B) \wedge (B \leftrightarrow C) \to
       (A \leftrightarrow C)$ --- тавтология.
\item \textbf{Отрицание:}
      $\neg(A \leftrightarrow B) \equiv A \leftrightarrow \neg B$.
\item $A \leftrightarrow B \equiv (A \wedge B) \vee (\neg A \wedge \neg B)$.
\end{enumerate}
\end{property}

%-----------------------------------------
\section{СДНФ и СКНФ}
%-----------------------------------------

\begin{definition}
\textbf{Литерал} --- переменная или её отрицание: $x_i$ или $\neg x_i$.

\textbf{Элементарная конъюнкция} --- конъюнкция литералов.

\textbf{Элементарная дизъюнкция} (клоз) --- дизъюнкция литералов.
\end{definition}

\begin{definition}
\textbf{Минтерм} для набора $(a_1,\ldots,a_n) \in \{0,1\}^n$ ---
элементарная конъюнкция $\displaystyle m_{a_1\cdots a_n} =
\bigwedge_{i=1}^n x_i^{a_i}$,
где $x_i^1 = x_i$ и $x_i^0 = \neg x_i$.
Минтерм равен 1 ровно на наборе $(a_1,\ldots,a_n)$.
\end{definition}

\begin{definition}
\textbf{СДНФ} (Совершенная Дизъюнктивная Нормальная Форма) функции $f$ ---
дизъюнкция всех минтермов, на которых $f = 1$:
\[
\text{СДНФ}(f) = \bigvee_{(a_1,\ldots,a_n):\; f(a_1,\ldots,a_n)=1}
m_{a_1\cdots a_n}.
\]
\end{definition}

\begin{definition}
\textbf{Макстерм} для набора $(a_1,\ldots,a_n)$ ---
элементарная дизъюнкция $\displaystyle M_{a_1\cdots a_n} =
\bigvee_{i=1}^n x_i^{1-a_i}$.
Макстерм равен 0 ровно на наборе $(a_1,\ldots,a_n)$.
\end{definition}

\begin{definition}
\textbf{СКНФ} (Совершенная Конъюнктивная Нормальная Форма) функции $f$ ---
конъюнкция всех макстермов, на которых $f = 0$:
\[
\mathrm{СКНФ}(f) = \bigwedge_{(a_1,\ldots,a_n):\; f(a_1,\ldots,a_n)=0}
M_{a_1\cdots a_n}.
\]
\end{definition}

\begin{example}
Пусть $f(A,B) = A \to B$.

\begin{center}
\begin{tabular}{|c|c|c|}
\hline
$A$ & $B$ & $A \to B$ \\ \hline
0 & 0 & 1 \\
0 & 1 & 1 \\
1 & 0 & 0 \\
1 & 1 & 1 \\ \hline
\end{tabular}
\end{center}

СДНФ: наборы с $f=1$: $(0,0),(0,1),(1,1)$.
\[
\text{СДНФ} = (\neg A \wedge \neg B) \vee (\neg A \wedge B)
\vee (A \wedge B).
\]

СКНФ: набор с $f=0$: $(1,0)$.
\[
\mathrm{СКНФ} = (\neg A \vee B).
\]
\end{example}

\begin{theorem}
Любая булева функция имеет единственные СДНФ и СКНФ (с точностью
до порядка слагаемых/множителей).
\end{theorem}
\begin{proof}
\textbf{Существование.} Построим СДНФ по определению:
для каждого набора $(a_1,\ldots,a_n)$, на котором $f=1$, образуем
минтерм. Их дизъюнкция совпадает с $f$, так как на каждом
наборе ровно один минтерм равен 1.

\textbf{Единственность.} Если бы существовала другая СДНФ,
в ней отсутствовал бы некоторый минтерм $m$ (на наборе которого $f=1$)
или присутствовал лишний. В первом случае СДНФ равна 0 на этом наборе,
во втором --- 1 на наборе, где $f=0$. Противоречие.
\end{proof}

%-----------------------------------------
\section{Выражение операций через базисные}
%-----------------------------------------

\begin{theorem}
Любая булева функция выражается через $\{\neg, \vee\}$,
через $\{\neg, \wedge\}$, а также через $\{\neg, \vee, \wedge\}$.
\end{theorem}

\begin{proof}
Достаточно выразить все операции через два выбранных базиса.

\medskip
\textbf{Базис $\{\neg, \vee\}$:}
\[
A \wedge B = \neg(\neg A \vee \neg B)
\quad\text{(де Морган)},
\]
\[
A \to B = \neg A \vee B,
\]
\[
A \leftrightarrow B = (\neg A \vee B) \wedge (\neg B \vee A)
= \neg(\neg(\neg A \vee B) \vee \neg(\neg B \vee A)).
\]

\medskip
\textbf{Базис $\{\neg, \wedge\}$:}
\[
A \vee B = \neg(\neg A \wedge \neg B)
\quad\text{(де Морган)},
\]
\[
A \to B = \neg A \vee B = \neg(\neg(\neg A) \wedge \neg B)
= \neg(A \wedge \neg B),
\]
\[
A \leftrightarrow B = (A \wedge B) \vee (\neg A \wedge \neg B)
= \neg\!\bigl(\neg(A \wedge B) \wedge \neg(\neg A \wedge \neg B)\bigr).
\]
\end{proof}

%-----------------------------------------
\section{Штрих Шеффера и его свойства}
%-----------------------------------------

\begin{definition}
\textbf{Штрих Шеффера} (операция NAND):
\[
A \mid B \stackrel{\mathrm{def}}{=} \neg(A \wedge B).
\]
Таблица истинности:
\begin{center}
\begin{tabular}{|c|c|c|}
\hline
$A$ & $B$ & $A \mid B$ \\ \hline
0 & 0 & 1 \\
0 & 1 & 1 \\
1 & 0 & 1 \\
1 & 1 & 0 \\ \hline
\end{tabular}
\end{center}
\end{definition}

\begin{theorem}
Множество $\{|\}$ является \emph{функционально полным} базисом:
любая булева функция выражается через штрих Шеффера.
\end{theorem}

\begin{proof}
Выразим $\neg$ и $\wedge$ через $|$:
\[
\neg A = A \mid A,
\]
\[
A \wedge B = \neg(A \mid B) = (A \mid B) \mid (A \mid B).
\]
Поскольку $\{\neg, \wedge\}$ --- полный базис, а обе операции
выражаются через $|$, то $\{|\}$ --- полный базис.
\end{proof}

\begin{property}[Свойства штриха Шеффера]
\begin{enumerate}[(i)]
\item \textbf{Коммутативность:} $A \mid B = B \mid A$.
\item \textbf{Не ассоциативен:} $(A \mid B) \mid C \neq A \mid (B \mid C)$
      в общем случае.
\item $A \mid A = \neg A$.
\item $A \mid 0 = 1$, $\quad A \mid 1 = \neg A$.
\end{enumerate}
\end{property}

%-----------------------------------------
\section{Стрелка Пирса и её свойства}
%-----------------------------------------

\begin{definition}
\textbf{Стрелка Пирса} (операция NOR):
\[
A \downarrow B \stackrel{\mathrm{def}}{=} \neg(A \vee B).
\]
Таблица истинности:
\begin{center}
\begin{tabular}{|c|c|c|}
\hline
$A$ & $B$ & $A \downarrow B$ \\ \hline
0 & 0 & 1 \\
0 & 1 & 0 \\
1 & 0 & 0 \\
1 & 1 & 0 \\ \hline
\end{tabular}
\end{center}
\end{definition}

\begin{theorem}
Множество $\{\downarrow\}$ является \emph{функционально полным} базисом.
\end{theorem}

\begin{proof}
Выразим $\neg$ и $\vee$ через $\downarrow$:
\[
\neg A = A \downarrow A,
\]
\[
A \vee B = \neg(A \downarrow B) = (A \downarrow B) \downarrow (A \downarrow B).
\]
Поскольку $\{\neg, \vee\}$ --- полный базис, то $\{\downarrow\}$ полон.
\end{proof}

\begin{property}[Свойства стрелки Пирса]
\begin{enumerate}[(i)]
\item \textbf{Коммутативность:} $A \downarrow B = B \downarrow A$.
\item \textbf{Не ассоциативна.}
\item $A \downarrow A = \neg A$.
\item $A \downarrow 0 = \neg A$, $\quad A \downarrow 1 = 0$.
\item Конъюнкция:
      $A \wedge B = (A \downarrow A) \downarrow (B \downarrow B)$.
\end{enumerate}
\end{property}

%-----------------------------------------
\section{Понятие множества. Способы задания. Операции.}
%-----------------------------------------

\begin{definition}
\textbf{Множество} --- совокупность различимых объектов (элементов),
рассматриваемая как единое целое. Обозначение: $A, B, \ldots$;
принадлежность: $x \in A$ или $x \notin A$.
\end{definition}

\subsection{Способы задания множества}

\begin{enumerate}
\item \textbf{Перечисление:} $A = \{1, 2, 3\}$.
\item \textbf{Характеристическое свойство:} $A = \{x \mid P(x)\}$,
      где $P(x)$ --- некоторое условие.
      \begin{example} $A = \{x \in \mathbb{Z} \mid x^2 < 10\}
      = \{-3,-2,-1,0,1,2,3\}$. \end{example}
\item \textbf{Диаграмма Эйлера--Венна} (графическое задание).
\end{enumerate}

\subsection{Операции на множествах}

\begin{definition}
Пусть $A, B \subseteq U$ (универсум).
\begin{itemize}
\item \textbf{Объединение:} $A \cup B = \{x \mid x \in A \text{ или } x \in B\}$.
\item \textbf{Пересечение:} $A \cap B = \{x \mid x \in A \text{ и } x \in B\}$.
\item \textbf{Разность:} $A \setminus B = \{x \mid x \in A \text{ и } x \notin B\}$.
\item \textbf{Симметрическая разность:}
      $A \triangle B = (A \setminus B) \cup (B \setminus A)$.
\item \textbf{Дополнение:} $\overline{A} = U \setminus A$.
\item \textbf{Включение:} $A \subseteq B \iff \forall x\,(x \in A \Rightarrow x \in B)$.
\end{itemize}
\end{definition}

\begin{property}[Законы де Моргана для множеств]
\[
\overline{A \cup B} = \overline{A} \cap \overline{B},
\qquad
\overline{A \cap B} = \overline{A} \cup \overline{B}.
\]
\end{property}

%-----------------------------------------
\section{Декартово произведение}
%-----------------------------------------

\begin{definition}
\textbf{Декартово произведение} множеств $A$ и $B$:
\[
A \times B = \{(a,b) \mid a \in A,\; b \in B\}.
\]
Элементы --- \emph{упорядоченные пары}.
\end{definition}

\begin{theorem}[Мощность декартова произведения]
$|A \times B| = |A| \cdot |B|$.
\end{theorem}

\begin{proof}
Для каждого из $|A|$ элементов $a \in A$ существует ровно $|B|$
пар $(a,b)$ при $b \in B$. Все пары различны, значит, их всего
$|A| \cdot |B|$.
\end{proof}

\begin{theorem}[Дистрибутивность декартова произведения]
\[
A \times (B \cup C) = (A \times B) \cup (A \times C),
\]
\[
A \times (B \cap C) = (A \times B) \cap (A \times C),
\]
\[
A \times (B \setminus C) = (A \times B) \setminus (A \times C).
\]
\end{theorem}

\begin{proof}
Докажем первое равенство. Пусть $(x,y) \in A \times (B \cup C)$.
Тогда $x \in A$ и $y \in B \cup C$, то есть $y \in B$ или $y \in C$.
Значит, $(x,y) \in A \times B$ или $(x,y) \in A \times C$,
т.\,е.\ $(x,y) \in (A\times B) \cup (A\times C)$.
Обратное включение аналогично.
\end{proof}

%-----------------------------------------
\section{Мощность. Формула включений и исключений.}
%-----------------------------------------

\begin{definition}
\textbf{Мощность} (кардинальное число) конечного множества $A$ ---
количество его элементов, обозначается $|A|$ или $\#A$.
\end{definition}

\begin{theorem}[Формула включений и исключений (для двух множеств)]\label{thm:incl-excl-2}
\[
|A \cup B| = |A| + |B| - |A \cap B|.
\]
\end{theorem}

\begin{proof}
Каждый элемент $x$ учтём в левой и правой частях:
\begin{itemize}
\item Если $x \in A \setminus B$: слева 1, справа $1 + 0 - 0 = 1$.
\item Если $x \in B \setminus A$: слева 1, справа $0 + 1 - 0 = 1$.
\item Если $x \in A \cap B$: слева 1, справа $1 + 1 - 1 = 1$.
\end{itemize}
\end{proof}

\begin{theorem}[Формула включений и исключений (для $n$ множеств)]
\[
\Bigl|\bigcup_{i=1}^n A_i\Bigr|
= \sum_{i}|A_i|
- \sum_{i<j}|A_i \cap A_j|
+ \sum_{i<j<k}|A_i \cap A_j \cap A_k|
- \cdots
+ (-1)^{n-1}|A_1 \cap \cdots \cap A_n|.
\]
\end{theorem}

\begin{proof}
Индукция по $n$. База $n=2$ доказана выше (см.~теорему~\ref{thm:incl-excl-2}).
Шаг: $\bigcup_{i=1}^n A_i = \bigl(\bigcup_{i=1}^{n-1} A_i\bigr) \cup A_n$.
Применяем формулу для двух множеств и предположение индукции, используя
$\bigl(\bigcup_{i=1}^{n-1} A_i\bigr) \cap A_n = \bigcup_{i=1}^{n-1}(A_i \cap A_n)$.
\end{proof}

%-----------------------------------------
\section{Диаграммы Эйлера и Эйлера--Венна}
%-----------------------------------------

\begin{definition}
\textbf{Диаграмма Эйлера} --- изображение множеств замкнутыми
кривыми (кругами). Изображаются \emph{только реально существующие}
пересечения.
\end{definition}

\begin{definition}
\textbf{Диаграмма Эйлера--Венна} --- изображение $n$ множеств,
при котором показаны \emph{все} $2^n$ возможных пересечений
(некоторые из них могут быть пустыми). Для трёх множеств ---
три круга, для четырёх --- эллипсы и т.\,д.
\end{definition}

\begin{note}
\textbf{Отличие}: диаграмма Эйлера отражает \emph{реальную} структуру
отношений между множествами; диаграмма Эйлера--Венна всегда
изображает \emph{все} возможные области, даже пустые.
Диаграмма Венна универсальна для записи формул с неизвестными
множествами.
\end{note}

%-----------------------------------------
\section{Счётные и несчётные множества}
%-----------------------------------------

\begin{definition}
Два множества $A$ и $B$ \textbf{равномощны} ($|A| = |B|$),
если существует биекция $f\colon A \to B$.
\end{definition}

\begin{definition}
Множество $A$ \textbf{счётно}, если оно равномощно натуральным
числам: $|A| = |\mathbb{N}|$. Обозначение: $|A| = \aleph_0$.

Множество называется \textbf{не более чем счётным}, если оно конечно
или счётно.
\end{definition}

\begin{definition}
Множество $A$ \textbf{несчётно}, если оно не является конечным
и не является счётным.
\end{definition}

\begin{example}
$\mathbb{Z}$, $\mathbb{Q}$, $\mathbb{N}\times\mathbb{N}$ --- счётны.
$\mathbb{R}$, интервал $(0,1)$, $2^{\mathbb{N}}$ --- несчётны.
\end{example}

\begin{theorem}[Кантор: $(0,1)$ несчётен]
Интервал $(0,1)$ несчётен.
\end{theorem}

\begin{proof}[Доказательство (диагональный метод Кантора)]
Предположим, что $(0,1)$ счётен: $x_1, x_2, x_3, \ldots$ ---
все его элементы. Запишем каждый в виде бесконечной десятичной
дроби: $x_i = 0.d_{i1}d_{i2}d_{i3}\cdots$

Построим число $y = 0.b_1 b_2 b_3 \cdots$, где
\[
b_i = \begin{cases} 1, & d_{ii} \neq 1 \\ 2, & d_{ii} = 1. \end{cases}
\]
Тогда $y \in (0,1)$, но $y \neq x_i$ для любого $i$
(т.\,к.\ $b_i \neq d_{ii}$). Противоречие с полнотой перечисления.
\end{proof}

%%%%%%%%%%%%%%%%%%%%%%%%%%%%%%%%%%%%%%%%%%%%%%%%%%%%%%%%%%%%%
\part{Теория графов}
%%%%%%%%%%%%%%%%%%%%%%%%%%%%%%%%%%%%%%%%%%%%%%%%%%%%%%%%%%%%%

\section{Основные понятия теории графов}

\begin{definition}
\textbf{Граф} $G = (V, E)$ --- пара, где $V$ --- непустое множество
\emph{вершин}, $E \subseteq \binom{V}{2}$ --- множество \emph{рёбер}
(неупорядоченных пар вершин). Порядок графа $|V| = n$,
размер $|E| = m$.
\end{definition}

\begin{definition}
\textbf{Псевдограф} --- граф, допускающий \emph{петли} (ребро
из вершины в себя) и \emph{кратные рёбра} (несколько рёбер между
одной парой вершин).
\end{definition}

\begin{definition}
\textbf{Полный граф} $K_n$ --- простой граф на $n$ вершинах,
в котором каждая пара вершин соединена ребром. Число рёбер:
$|E(K_n)| = \binom{n}{2} = \dfrac{n(n-1)}{2}$.
\end{definition}

\begin{definition}
\textbf{Двудольный граф} $G = (V_1 \cup V_2, E)$ --- граф,
вершины которого можно разбить на два множества $V_1, V_2$
так, что каждое ребро соединяет вершину из $V_1$ с вершиной
из $V_2$. \textbf{Полный двудольный граф} $K_{m,n}$ содержит
все $m \cdot n$ таких рёбер.
\end{definition}

\begin{definition}
\textbf{Колесо} $W_n$ ($n \geq 3$) --- граф, полученный добавлением
одной вершины («ступицы») к циклу $C_n$ и соединением её со всеми
вершинами цикла. $|V(W_n)| = n+1$, $|E(W_n)| = 2n$.
\end{definition}

\begin{definition}
\textbf{Цепь} (walk) длины $k$ --- последовательность
$v_0 e_1 v_1 e_2 \cdots e_k v_k$, где $e_i = \{v_{i-1},v_i\}$.
\textbf{Простой путь} (path) --- цепь без повторяющихся вершин.
\textbf{Цикл} (cycle) --- замкнутая простая цепь ($v_0 = v_k$,
все промежуточные вершины различны).
\end{definition}

\begin{definition}
\textbf{Объединение} графов $G_1 = (V_1,E_1)$ и $G_2=(V_2,E_2)$
(предполагаем $V_1 \cap V_2 = \varnothing$):
\[
G_1 \cup G_2 = (V_1 \cup V_2,\; E_1 \cup E_2).
\]
\textbf{Соединение} $G_1 + G_2$ --- объединение плюс все рёбра
между $V_1$ и $V_2$.
\end{definition}

\begin{definition}
Графы $G_1 = (V_1,E_1)$ и $G_2 = (V_2,E_2)$ \textbf{изоморфны}
($G_1 \cong G_2$), если существует биекция $\varphi\colon V_1 \to V_2$
такая, что
\[
\{u,v\} \in E_1 \iff \{\varphi(u),\varphi(v)\} \in E_2.
\]
\end{definition}

%-----------------------------------------
\section{Маршрут, путь, связность}
%-----------------------------------------

\begin{definition}
\textbf{Маршрут} (walk) от $u$ до $v$ --- последовательность
вершин $u = v_0, v_1, \ldots, v_k = v$, где каждые две соседние
вершины смежны. Длина маршрута --- число рёбер $k$.
Маршрут называется \textbf{замкнутым}, если $u = v$.

\textbf{Путь} (path) --- маршрут без повторяющихся вершин.

\textbf{Граф $G$ связен}, если между любыми двумя его вершинами
существует маршрут (а значит, и путь).
\end{definition}

\begin{theorem}[О компонентах связности]
Отношение «существует маршрут от $u$ до $v$» является
\emph{отношением эквивалентности} на $V(G)$.
Классы эквивалентности называются \textbf{компонентами связности}.
\end{theorem}

\begin{proof}
\begin{itemize}
\item \textbf{Рефлексивность:} маршрут нулевой длины из $v$ в $v$.
\item \textbf{Симметричность:} если $v_0,\ldots,v_k$ --- маршрут из $u$ в $v$,
      то $v_k,\ldots,v_0$ --- маршрут из $v$ в $u$.
\item \textbf{Транзитивность:} конкатенация двух маршрутов.
\end{itemize}
\end{proof}

%-----------------------------------------
\section{Критерий связности}
%-----------------------------------------

\begin{definition}
Обозначим через $c(G)$ число компонент связности графа $G$.
\end{definition}

\begin{theorem}[Критерий связности]
Граф $G$ на $n$ вершинах и $m$ рёбрах \emph{связен} тогда и только
тогда, когда $m \geq n - 1$ и граф не имеет изолированных вершин ---
нет, более точно:

Граф связен тогда и только тогда, когда $c(G) = 1$.
\end{theorem}

\begin{theorem}[Необходимое условие]
Если граф $G$ имеет $n$ вершин, $m$ рёбер и $k$ компонент связности,
то $m \geq n - k$.
\end{theorem}

\begin{proof}
В каждой компоненте $n_i$ вершин и $m_i$ рёбер.
Дерево на $n_i$ вершинах имеет $n_i - 1$ рёбер, а компонента
содержит хотя бы столько: $m_i \geq n_i - 1$.
Суммируя:
\[
m = \sum_{i=1}^k m_i \geq \sum_{i=1}^k (n_i - 1) = n - k.
\]
\end{proof}

%-----------------------------------------
\section{Расстояние, диаметр, центр, обхват}
%-----------------------------------------

\begin{definition}
В связном графе $G$ \textbf{расстояние} $d(u,v)$ --- длина
кратчайшего пути между $u$ и $v$.

\textbf{Эксцентриситет} вершины $v$: $e(v) = \max_{u \in V} d(u,v)$.

\textbf{Диаметр:} $\mathrm{diam}(G) = \max_{u,v} d(u,v) = \max_v e(v)$.

\textbf{Радиус:} $\mathrm{rad}(G) = \min_v e(v)$.

\textbf{Центр} --- множество вершин с минимальным эксцентриситетом:
$Z(G) = \{v \mid e(v) = \mathrm{rad}(G)\}$.

\textbf{Обхват} $g(G)$ --- длина кратчайшего цикла в $G$
(для ациклических графов $g(G) = \infty$).
\end{definition}

\begin{example}
Для цикла $C_n$: $\mathrm{diam}(C_n) = \lfloor n/2 \rfloor$,
$\mathrm{rad}(C_n) = \lfloor n/2 \rfloor$, $g(C_n) = n$.
Для $K_n$ ($n \geq 2$): $\mathrm{diam} = 1$, $g = 3$ (при $n \geq 3$).
\end{example}

%-----------------------------------------
\section{Эйлеров путь и цикл}
%-----------------------------------------

\begin{definition}
\textbf{Эйлеров путь} --- маршрут, проходящий по каждому ребру
ровно один раз.
\textbf{Эйлеров цикл} --- замкнутый эйлеров путь.
Граф называется \textbf{эйлеровым}, если в нём есть эйлеров цикл.
\end{definition}

\begin{theorem}[Эйлер, 1736]
Связный граф $G$ \emph{имеет эйлеров цикл} тогда и только тогда,
когда степень каждой его вершины чётна.
\end{theorem}

\begin{proof}
\textbf{Необходимость.} При обходе эйлерова цикла каждый
раз, когда маршрут входит в вершину $v$, он из неё выходит.
Поэтому каждое вхождение даёт вклад 2 в степень $v$. Итого $\deg v$ чётна.

\textbf{Достаточность.} Применяем алгоритм:
\begin{enumerate}
\item Начинаем из любой вершины, идём по рёбрам, не повторяя их,
      пока не вернёмся (получаем замкнутый маршрут $W$).
      Это возможно, так как все степени чётны.
\item Если $W$ не охватывает всех рёбер, то среди непройденных рёбер
      есть связная компонента, смежная с $W$ (граф связен).
      Начинаем новый замкнутый маршрут из вершины смежности и вставляем
      его в $W$.
\item Повторяем до использования всех рёбер.
\end{enumerate}
\end{proof}

\begin{theorem}
Связный граф имеет \emph{эйлеров путь} (не цикл) тогда и только тогда,
когда ровно две вершины имеют нечётную степень.
\end{theorem}

\begin{proof}
Добавим ребро между двумя нечётными вершинами $u, v$. Тогда
все степени станут чётными, граф эйлеров --- существует эйлеров цикл.
Удалив добавленное ребро, получаем эйлеров путь из $u$ в $v$.
\end{proof}

%-----------------------------------------
\section{Гамильтонов путь. Теорема Дирака.}\label{sec:dirac}
%-----------------------------------------

\begin{definition}
\textbf{Гамильтонов путь} --- простой путь, проходящий через
каждую вершину ровно один раз.
\textbf{Гамильтонов цикл} --- замкнутый гамильтонов путь.
Граф называется \textbf{гамильтоновым}, если в нём есть гамильтонов цикл.
\end{definition}

\begin{theorem}[Дирак, 1952 --- гамильтонов цикл]\label{thm:dirac}
Пусть $G$ --- простой граф на $n \geq 3$ вершинах, и для каждой
вершины $v$: $\deg v \geq \dfrac{n}{2}$. Тогда $G$ гамильтонов.
\end{theorem}

\begin{proof}
\textbf{Шаг 1. Связность.} Возьмём два несмежных $u,v$.
$|\Gamma(u)| \geq n/2$, $|\Gamma(v)| \geq n/2$, $|\Gamma(u)\cup\Gamma(v)| \leq n$.
По формуле включений-исключений $|\Gamma(u)\cap\Gamma(v)| \geq n - n = 0$... 
нет, точнее: $|\Gamma(u)| + |\Gamma(v)| \geq n$, а $|\Gamma(u) \cup \Gamma(v)| \leq n$,
значит $|\Gamma(u) \cap \Gamma(v)| \geq 0$. Из $|\Gamma(u)| + |\Gamma(v)| \geq n$
и $|\Gamma(u) \cup \Gamma(v) \cup \{u,v\}| \leq n$ получаем общего соседа.
Граф связен.

\textbf{Шаг 2. Длиннейший путь.} Пусть $P = v_1 v_2 \cdots v_k$ ---
простой путь максимальной длины. Если $k < n$, то все соседи
$v_1$ и $v_k$ лежат на $P$.

\textbf{Шаг 3. Цикл.} Если $v_1 v_k \in E$, есть цикл длины $k$.
Иначе: пусть $v_1$ смежна с $v_j$ и $v_k$ смежна с $v_{j-1}$.
Число таких индексов $j$: $|\Gamma(v_k) \cap \{v_1,\ldots,v_{k-1}\}| \geq n/2$.
Число соседей $v_1$ на $P$: $\geq n/2$. Поскольку $k \leq n$,
$\deg(v_1) + \deg(v_k) \geq n > k-1$, значит существует $j$ такой,
что $v_j \in \Gamma(v_1)$ и $v_{j-1} \in \Gamma(v_k)$ одновременно.
Цикл: $v_1 v_2 \cdots v_{j-1} v_k v_{k-1} \cdots v_j v_1$.

\textbf{Шаг 4.} Если цикл имеет длину $k < n$, из связности
есть вершина вне цикла, смежная с ним, что даёт более длинный путь ---
противоречие. Значит $k = n$.
\end{proof}

%-----------------------------------------
\section{Теорема Дирака про гамильтонов цикл}
%-----------------------------------------

(Уже доказана в \S\ref{sec:dirac} --- см.~теорему~\ref{thm:dirac}.)

%-----------------------------------------
\section{Деревья}
%-----------------------------------------

\begin{definition}
\textbf{Дерево} --- связный ациклический граф.
\textbf{Лес} --- ациклический граф (объединение деревьев).
\textbf{Остовное дерево} графа $G$ --- подграф $T \subseteq G$,
являющийся деревом и содержащий все вершины $G$.
\textbf{Висячая вершина} (лист) --- вершина степени 1.
\end{definition}

\begin{theorem}[Эквивалентные определения дерева]
Для простого графа $G$ на $n$ вершинах следующие утверждения
эквивалентны:
\begin{enumerate}[(1)]
\item $G$ --- дерево (связен и ацикличен).
\item $G$ связен и $|E| = n-1$.
\item $G$ ацикличен и $|E| = n-1$.
\item Между любыми двумя вершинами существует единственный простой путь.
\item $G$ связен, но удаление любого ребра делает его несвязным.
\end{enumerate}
\end{theorem}

\begin{proof}
$(1) \Rightarrow (2)$: Индукция по $n$.
База $n=1$: $|E|=0=1-1$. Шаг: дерево на $n \geq 2$ вершинах имеет лист $v$ (см.~теорему~\ref{thm:leaves} ниже). Удалим $v$ и инцидентное ему ребро --- получим дерево на $n-1$ вершинах с $n-2$ рёбрами. Итого $n-1$ ребро.

$(2) \Rightarrow (3)$: Если граф связен и $|E|=n-1$, то он ацикличен.
(Если бы был цикл, удаление ребра из цикла сохраняло бы связность,
а граф на $n$ вершинах с $< n-1$ рёбрами не может быть связным.)

$(3) \Rightarrow (1)$: Если ацикличен и $|E|=n-1$,
то в каждой из $k$ компонент связности --- дерево (ацикличное связное),
т.\,е.\ $n_i - 1$ рёбер. $|E| = \sum(n_i - 1) = n - k = n-1$,
значит $k = 1$.

$(1) \Leftrightarrow (4)$: Связность гарантирует существование пути.
Ацикличность гарантирует единственность.
\end{proof}

\begin{theorem}[Наличие висячей вершины]\label{thm:leaves}
Каждое дерево на $n \geq 2$ вершинах имеет хотя бы одну висячую вершину.
\end{theorem}

\begin{proof}
$\sum_{v}\deg v = 2|E| = 2(n-1)$.
Если бы все вершины имели $\deg \geq 2$, то
$\sum \deg \geq 2n > 2(n-1)$ --- противоречие.
\end{proof}

%-----------------------------------------
\section{Планарность. Теорема Эйлера. Следствия.}
%-----------------------------------------

\begin{definition}
Граф \textbf{планарен}, если его можно нарисовать на плоскости
без пересечений рёбер. Такое рисование --- \textbf{плоская укладка}.
\textbf{Грань} плоской укладки --- связная область плоскости,
ограниченная рёбрами (включая неограниченную внешнюю грань).
\end{definition}

\begin{theorem}[Формула Эйлера для плоских графов]
Пусть $G$ --- связный плоский граф с $n$ вершинами, $m$ рёбрами
и $f$ гранями. Тогда:
\[
n - m + f = 2.
\]
\end{theorem}

\begin{proof}
Индукция по числу рёбер $m$.

\textbf{База.} $m = n-1$ (дерево): $f = 1$ (только внешняя грань),
$n - (n-1) + 1 = 2$. \checkmark

\textbf{Шаг.} Пусть $G$ не является деревом (есть цикл).
Возьмём ребро $e$ на цикле. Граф $G' = G - e$ связен,
имеет $n$ вершин, $m-1$ рёбер, $f' = f-1$ граней
(удаление ребра из цикла объединяет две грани в одну).
По предположению индукции: $n - (m-1) + (f-1) = 2$,
откуда $n - m + f = 2$.
\end{proof}

\begin{corollary}\label{cor:euler-edges}
Если $G$ --- простой связный планарный граф с $n \geq 3$ вершинами
и $m$ рёбрами, то $m \leq 3n - 6$.
\end{corollary}

\begin{proof}
Каждая грань ограничена $\geq 3$ рёбрами. Каждое ребро граничит
с $\leq 2$ гранями. Значит $3f \leq 2m$, т.\,е.\ $f \leq 2m/3$.
Из формулы Эйлера: $f = 2 - n + m$, откуда
$2 - n + m \leq 2m/3$, $m/3 \leq n-2$, $m \leq 3n-6$.
\end{proof}

\begin{corollary}
Если $G$ --- простой связный планарный граф без треугольников
($g \geq 4$), то $m \leq 2n - 4$.
\end{corollary}

\begin{proof}
Каждая грань ограничена $\geq 4$ рёбрами: $4f \leq 2m$, $f \leq m/2$.
$2 - n + m \leq m/2$, $m/2 \leq n-2$, $m \leq 2n-4$.
\end{proof}

\begin{corollary}
$K_5$ и $K_{3,3}$ не планарны.
\end{corollary}

\begin{proof}
$K_5$: $n=5$, $m=10$, $3n-6=9 < 10$. Противоречие.
$K_{3,3}$: $n=6$, $m=9$, $g=4$, $2n-4=8 < 9$. Противоречие.
\end{proof}

%-----------------------------------------
\section{Теорема Куратовского}
%-----------------------------------------

\begin{definition}
\textbf{Подразделение} графа $H$ --- граф, полученный заменой
некоторых рёбер $H$ простыми путями (вставкой вершин степени 2
на рёбра). Если $H' $ является подразделением $H$, говорят,
что $H'$ --- \emph{подразделение} $H$.

Граф $G$ содержит \textbf{подграф, гомеоморфный} $H$,
если в $G$ есть подразделение $H$.
\end{definition}

\begin{theorem}[Куратовский, 1930]
Граф $G$ \emph{планарен} тогда и только тогда, когда он не содержит
подграфа, гомеоморфного $K_5$ или $K_{3,3}$.
\end{theorem}

\begin{proof}[Идея доказательства]
\textbf{Необходимость.} Если $G$ планарен, любой его подграф
тоже планарен. Но подразделения $K_5$ и $K_{3,3}$ не планарны
(они имеют плоские укладки эквивалентные $K_5$/$K_{3,3}$, которые
не планарны по следствиям теоремы Эйлера).

\textbf{Достаточность.} (Сложная часть.) Индукция по числу рёбер.
Если $G$ не планарен, то существует минимальный непланарный подграф $H$.
Анализ структуры $H$ (связность, степени вершин, удаление вершин
и рёбер) показывает, что $H$ должен содержать подразделение
$K_5$ или $K_{3,3}$.
(Полное доказательство требует рассмотрения нескольких случаев и
здесь опускается.)
\end{proof}

%-----------------------------------------
\section{Двойственный псевдограф}
%-----------------------------------------

\begin{definition}
Пусть $G$ --- связный плоский граф с укладкой. \textbf{Двойственный
псевдограф} $G^*$ строится следующим образом:
\begin{itemize}
\item Каждой грани $f_i$ графа $G$ соответствует вершина $f_i^*$ в $G^*$.
\item Каждому ребру $e$ графа $G$, разделяющему грани $f_i$ и $f_j$,
      соответствует ребро $\{f_i^*, f_j^*\}$ в $G^*$.
\end{itemize}
Если $e$ граничит с одной гранью (мост), то в $G^*$ появляется петля.
\end{definition}

\begin{theorem}
Двойственный граф $G^{**} \cong G$ (для плоского графа без мостов).
Также: $|V(G^*)| = f$, $|E(G^*)| = m$, $|E_{\text{faces}}(G^*)| = n$,
где $f, m, n$ --- число граней, рёбер и вершин $G$.
\end{theorem}

\begin{proof}
По построению: вершины $G^*$ --- грани $G$ ($f$ штук),
рёбра $G^*$ --- рёбра $G$ ($m$ штук), грани $G^*$ ---
вершины $G$ ($n$ штук). Формула Эйлера для $G^*$:
$f - m + n = 2$, что совпадает с формулой для $G$.
\end{proof}

%-----------------------------------------
\section{Раскраски}
%-----------------------------------------

\begin{definition}
\textbf{Правильная раскраска} вершин графа в $k$ цветов ---
отображение $c\colon V \to \{1,\ldots,k\}$ такое, что
$c(u) \neq c(v)$ для каждого $\{u,v\} \in E$.
\textbf{Хроматическое число} $\chi(G)$ --- минимальное $k$,
при котором $G$ имеет правильную $k$-раскраску.
\end{definition}

\begin{theorem}[$(\Delta+1)$-раскраска; вершинная раскраска]
Для любого графа $G$: $\chi(G) \leq \Delta(G) + 1$,
где $\Delta(G)$ --- максимальная степень вершины.
\end{theorem}

\begin{proof}
Жадный алгоритм: упорядочиваем вершины $v_1,\ldots,v_n$.
Для каждой $v_i$ выбираем наименьший цвет, не использованный
соседями. Сосед $v_i$ среди $v_1,\ldots,v_{i-1}$ не более $\Delta$,
значит занято не более $\Delta$ цветов, свободен хотя бы один
из $\Delta+1$.
\end{proof}

\begin{theorem}[О 6 красках для планарных графов]
Каждый планарный граф можно правильно раскрасить в 6 цветов:
$\chi(G) \leq 6$.
\end{theorem}

\begin{proof}
Из $m \leq 3n-6$ (см.~следствие~\ref{cor:euler-edges}) следует, что существует вершина $v$ с $\deg v \leq 5$.
Индукция по $n$:
\begin{enumerate}
\item База $n \leq 6$: тривиально, хватает 6 цветов.
\item Шаг: есть вершина $v$ с $\deg v \leq 5$.
      По индукции $G - v$ раскрашивается в 6 цветов.
      Соседей $v$ не более 5, они заняли не более 5 цветов.
      Для $v$ свободен хотя бы 1 из 6.
\end{enumerate}
\end{proof}

%-----------------------------------------
\section{Теорема о 5 красках}
%-----------------------------------------

\begin{theorem}[Пять красок Кемпе--Хивуда]
Каждый планарный граф можно правильно раскрасить в 5 цветов.
\end{theorem}

\begin{proof}
Индукция по $n = |V(G)|$.

\textbf{База} $n \leq 5$: очевидно.

\textbf{Шаг.} Выберем вершину $v$ с $\deg v \leq 5$ (существует из следствия~\ref{cor:euler-edges}, так как $m \leq 3n-6 \Rightarrow \bar{d} < 6$). По индукции $G' = G - v$ раскрашивается в 5 цветов.

Если соседи $v$ используют $\leq 4$ цветов, раскрашиваем $v$
в 5-й. Готово.

Если соседи $v$ --- $v_1,\ldots,v_5$ --- используют все 5 цветов
(по кругу цвета $1,2,3,4,5$):

Рассмотрим подграф $H_{13}$, индуцированный вершинами цветов 1 и 3.
\begin{itemize}
\item Если $v_1$ и $v_3$ лежат в \emph{разных} компонентах $H_{13}$:
      в компоненте $v_1$ меняем цвета 1 и 3 местами. Теперь $v_3$
      и $v_1$ одного цвета 3, красим $v$ в цвет 1.
\item Если в одной компоненте $H_{13}$: существует цепь Кемпе между
      $v_1$ (цвет 1) и $v_3$ (цвет 3). Эта цепь «отделяет»
      $v_2$ от $v_4$ на плоскости. Применяем ту же процедуру к $H_{24}$
      --- $v_2$ и $v_4$ в разных компонентах, меняем цвета в компоненте
      $v_2$. Теперь $v_2$ получает цвет 4, и цвет 2 освобождается
      для вершины $v$.
\end{itemize}
\end{proof}

%-----------------------------------------
\section{Ориентированные графы}
%-----------------------------------------

\begin{definition}
\textbf{Ориентированный граф} (орграф) $D = (V, A)$ --- пара,
где $A \subseteq V \times V$ --- множество \textbf{дуг}
(упорядоченных пар вершин).

\textbf{Полустепень захода} $\deg^-(v)$ --- число дуг, входящих в $v$.
\textbf{Полустепень исхода} $\deg^+(v)$ --- число дуг, исходящих из $v$.
\[
\sum_{v \in V} \deg^+(v) = \sum_{v \in V} \deg^-(v) = |A|.
\]

\textbf{Сильно связный орграф} --- для любых $u,v$ существует
ориентированный маршрут из $u$ в $v$ и из $v$ в $u$.

\textbf{Эйлеров контур} в орграфе --- замкнутый маршрут,
проходящий по каждой дуге ровно один раз.
\end{definition}

\begin{theorem}[Эйлеров цикл в орграфе]
Сильно связный орграф $D$ имеет эйлеров контур тогда и только тогда,
когда для каждой вершины $v$:
$\deg^+(v) = \deg^-(v)$.
\end{theorem}

\begin{proof}
\textbf{Необходимость.} При каждом прохождении через $v$
маршрут использует одну входящую и одну исходящую дугу.
Значит, $\deg^+(v) = \deg^-(v)$.

\textbf{Достаточность.} Аналогично неориентированному случаю:
строим замкнутые маршруты и объединяем их, используя условие
$\deg^+ = \deg^-$ для обеспечения возможности продолжения.
\end{proof}

%-----------------------------------------
\section{Количество маркированных деревьев. Остовных деревьев.}
%-----------------------------------------

\begin{theorem}[Кэли, 1889]
Число помеченных деревьев (маркированных деревьев) на $n$ вершинах
равно $n^{n-2}$.
\end{theorem}

\begin{proof}[Доказательство через код Прюфера]
\textbf{Код Прюфера.} Дано маркированное дерево $T$ на вершинах
$\{1,\ldots,n\}$. Алгоритм построения кода:
\begin{enumerate}
\item Пока в дереве $> 2$ вершин:
      найти наименьший лист $l$, записать его \emph{соседа} $a_i$
      в код, удалить $l$ из дерева.
\end{enumerate}
Код $P(T) = (a_1,\ldots,a_{n-2}) \in \{1,\ldots,n\}^{n-2}$.

\textbf{Взаимно однозначное соответствие.}
Код однозначно определяет дерево:
на $i$-м шаге наименьший лист --- это наименьшая вершина,
не встречающаяся в $a_i,a_{i+1},\ldots,a_{n-2}$.

Таким образом, деревья биективны с последовательностями
$\{1,\ldots,n\}^{n-2}$, которых ровно $n^{n-2}$.
\end{proof}

\begin{theorem}[Кирхгоф; матрично-дерево теорема]
Число остовных деревьев связного графа $G$ равно любому
кофактору матрицы Лапласа $L(G) = D - A$, где $D$ ---
диагональная матрица степеней, $A$ --- матрица смежности.
\end{theorem}

\begin{proof}[Идея]
Матрица Лапласа: $L_{ij} = \begin{cases} \deg(i), & i=j \\ -1, & \{i,j\}\in E \\ 0, & \text{иначе} \end{cases}$.

По теореме Матрично-дерево Кирхгофа любой кофактор $C_{ii}$
(определитель матрицы $L$ с удалённой $i$-й строкой и $i$-м столбцом)
равен числу остовных деревьев. Доказательство использует формулу
Бине--Коши и разложение по деревьям.
\end{proof}

%-----------------------------------------
\section{Теорема Оре --- гамильтонов цикл}
%-----------------------------------------

\begin{theorem}[Оре, 1960]
Пусть $G$ --- простой граф на $n \geq 3$ вершинах, и для любых двух
\emph{несмежных} вершин $u, v$:
\[
\deg u + \deg v \geq n.
\]
Тогда $G$ гамильтонов.
\end{theorem}

\begin{proof}
\textbf{Шаг 1.} Предположим, что $G$ не гамильтонов.
Добавляем рёбра, не нарушая условия Оре, до тех пор, пока
граф не станет гамильтоновым (или пока нельзя добавить рёбро).
Получаем граф $H$ такой, что $H$ не гамильтонов, но $H + e$
гамильтонов для любого ребра $e \notin E(H)$.

\textbf{Шаг 2.} Пусть $\{u,v\} \notin E(H)$, тогда $H+\{u,v\}$
гамильтонов. Гамильтонов цикл в $H+\{u,v\}$ использует ребро
$\{u,v\}$ (иначе цикл был бы в $H$). Значит, существует
гамильтонов путь $v = w_1, w_2, \ldots, w_n = u$ в $H$.

\textbf{Шаг 3.} Пусть $S = \{i \mid \{u, w_i\} \in E(H)\}$,
$T = \{i \mid \{v, w_{i+1}\} \in E(H)\}$.
$|S| = \deg_H(u)$, $|T| = \deg_H(v)$, $S, T \subseteq \{1,\ldots,n-1\}$.
$|S| + |T| = \deg(u)+\deg(v) \geq n > n-1$, значит
$S \cap T \neq \varnothing$. Пусть $i \in S \cap T$:
\[
u \in \Gamma(w_i), \quad v \in \Gamma(w_{i+1}).
\]
Тогда $w_1 w_2 \cdots w_i u w_{i-1} \cdots w_{i+1} v$ ---
тоже... нет, строим цикл:
$w_1 \to w_2 \to \cdots \to w_i \to w_n \to w_{n-1} \to \cdots \to w_{i+1} \to w_1$
(используя $\{w_i, w_n\}=\{w_i,u\}$ и $\{w_{i+1},w_1\}=\{w_{i+1},v\}$).
Получаем гамильтонов цикл в $H$ --- противоречие.
\end{proof}

%-----------------------------------------
\section{Теорема Фари --- выпрямление графа}
%-----------------------------------------

\begin{theorem}[Фари, 1948]
Каждый простой планарный граф имеет прямолинейную плоскую укладку
(каждое ребро изображается отрезком прямой).
\end{theorem}

\begin{proof}
Индукция по числу вершин $n$.

\textbf{База} $n \leq 3$: очевидно.

\textbf{Шаг.} Пусть $G$ --- планарный граф без кратных рёбер.
Можно считать $G$ триангулированным (добавим рёбра в грани).
Выберем вершину $v$ с $\deg v \leq 5$ (существует по следствию~\ref{cor:euler-edges}).
Удалим $v$; по индукции $G - v$ имеет прямолинейную укладку.
Соседи $v$: $v_1,\ldots,v_k$ ($k \leq 5$) образуют многоугольник
на плоскости. Поместим $v$ внутри этого многоугольника
(можно сделать это корректно для $k \leq 5$).
Прямолинейные рёбра от $v$ к соседям не пересекают остальных.
\end{proof}

%%%%%%%%%%%%%%%%%%%%%%%%%%%%%%%%%%%%%%%%%%%%%%%%%%%%%%%%%%%%%
\part{Тригонометрия}
%%%%%%%%%%%%%%%%%%%%%%%%%%%%%%%%%%%%%%%%%%%%%%%%%%%%%%%%%%%%%

\section{Определения синуса и косинуса. ОТТ. Тригонометрическая окружность.}

\begin{definition}
\textbf{Тригонометрическая (единичная) окружность} --- окружность
радиуса 1 с центром в начале координат:
$x^2 + y^2 = 1$.

Угол $\alpha$ отсчитывается от положительного направления оси $Ox$
против часовой стрелки (положительное направление).

Точка на единичной окружности, соответствующая углу $\alpha$:
$(\cos\alpha, \sin\alpha)$.
\end{definition}

\begin{definition}
$\cos\alpha$ --- абсцисса точки единичной окружности,
соответствующей углу $\alpha$.

$\sin\alpha$ --- ордината той же точки.

$\tan\alpha = \dfrac{\sin\alpha}{\cos\alpha}$ ($\cos\alpha \neq 0$).

$\cot\alpha = \dfrac{\cos\alpha}{\sin\alpha}$ ($\sin\alpha \neq 0$).
\end{definition}

\begin{theorem}[Основное тригонометрическое тождество (ОТТ)]
\[
\sin^2\alpha + \cos^2\alpha = 1.
\]
\end{theorem}

\begin{proof}
Следует из определения: $(\cos\alpha)^2 + (\sin\alpha)^2$ ---
сумма квадратов координат точки на единичной окружности $= 1^2 = 1$.
\end{proof}

\begin{property}[Знаки в четвертях]
\begin{center}
\begin{tabular}{|c|c|c|}
\hline
Четверть & $\sin$ & $\cos$ \\ \hline
I ($0 < \alpha < \pi/2$) & $+$ & $+$ \\
II ($\pi/2 < \alpha < \pi$) & $+$ & $-$ \\
III ($\pi < \alpha < 3\pi/2$) & $-$ & $-$ \\
IV ($3\pi/2 < \alpha < 2\pi$) & $-$ & $+$ \\ \hline
\end{tabular}
\end{center}
\end{property}

%-----------------------------------------
\section{Формулы приведения}
%-----------------------------------------

Формулы приведения выражают тригонометрические функции углов
$\dfrac{\pi}{2} \pm \alpha$, $\pi \pm \alpha$, $\dfrac{3\pi}{2} \pm \alpha$,
$2\pi \pm \alpha$ через функции угла $\alpha$.

\textbf{Правило.} При смещении на $\pi/2$ или $3\pi/2$ функция меняется
на ко-функцию (sin $\leftrightarrow$ cos); при смещении на $\pi$ или $2\pi$
--- не меняется. Знак определяется по четверти, в которой лежит
исходный угол (считая $\alpha$ острым).

\begin{center}
\begin{tabular}{|l|l|l|}
\hline
Формула & $\sin$ & $\cos$ \\ \hline
$\pi/2 - \alpha$ & $\cos\alpha$ & $\sin\alpha$ \\
$\pi/2 + \alpha$ & $\cos\alpha$ & $-\sin\alpha$ \\
$\pi - \alpha$ & $\sin\alpha$ & $-\cos\alpha$ \\
$\pi + \alpha$ & $-\sin\alpha$ & $-\cos\alpha$ \\
$3\pi/2 - \alpha$ & $-\cos\alpha$ & $-\sin\alpha$ \\
$3\pi/2 + \alpha$ & $-\cos\alpha$ & $\sin\alpha$ \\
$2\pi - \alpha$ & $-\sin\alpha$ & $\cos\alpha$ \\
$2\pi + \alpha$ & $\sin\alpha$ & $\cos\alpha$ \\ \hline
\end{tabular}
\end{center}

\begin{proof}[Вывод формулы $\sin(\pi - \alpha) = \sin\alpha$]
На единичной окружности точка угла $\pi - \alpha$ имеет координаты
$(\cos(\pi-\alpha), \sin(\pi-\alpha))$. Отражение точки $(\cos\alpha, \sin\alpha)$
относительно оси $Oy$ даёт $(-\cos\alpha, \sin\alpha)$.
Значит $\cos(\pi-\alpha) = -\cos\alpha$, $\sin(\pi-\alpha) = \sin\alpha$.
\end{proof}

%-----------------------------------------
\section{Синус, косинус, тангенс суммы и разности}
%-----------------------------------------

\begin{theorem}[Формулы сложения]
\[
\sin(\alpha \pm \beta) = \sin\alpha\cos\beta \pm \cos\alpha\sin\beta,
\]
\[
\cos(\alpha \pm \beta) = \cos\alpha\cos\beta \mp \sin\alpha\sin\beta,
\]
\[
\tan(\alpha \pm \beta) = \frac{\tan\alpha \pm \tan\beta}
{1 \mp \tan\alpha\tan\beta}.
\]
\end{theorem}

\begin{proof}[Доказательство через координаты]
Рассмотрим на единичной окружности точки
$A = (\cos\alpha, \sin\alpha)$ и $B = (\cos\beta, \sin\beta)$.
Расстояние $|AB|^2 = (\cos\alpha - \cos\beta)^2 + (\sin\alpha - \sin\beta)^2
= 2 - 2\cos(\alpha - \beta)$.

С другой стороны, рассмотрим точки $C = (\cos(\alpha-\beta), \sin(\alpha-\beta))$
и $D = (1, 0)$:
$|CD|^2 = (\cos(\alpha-\beta)-1)^2 + \sin^2(\alpha-\beta)
= 2 - 2\cos(\alpha-\beta)$.
$|AB| = |CD|$ (хорды равного угла на единичной окружности).
Вычисляя напрямую:
\[
(\cos\alpha - \cos\beta)^2 + (\sin\alpha - \sin\beta)^2
= (\cos(\alpha-\beta)-1)^2 + \sin^2(\alpha-\beta),
\]
\[
2 - 2(\cos\alpha\cos\beta + \sin\alpha\sin\beta)
= 2 - 2\cos(\alpha-\beta).
\]
Значит $\cos(\alpha-\beta) = \cos\alpha\cos\beta + \sin\alpha\sin\beta$.

Заменяем $\beta \to -\beta$, используя чётность/нечётность:
$\cos(\alpha+\beta) = \cos\alpha\cos\beta - \sin\alpha\sin\beta$.

Для синуса: $\sin(\alpha+\beta) = \cos(\pi/2 - (\alpha+\beta))
= \cos((\pi/2 - \alpha) - \beta)$
$= \cos(\pi/2-\alpha)\cos\beta + \sin(\pi/2-\alpha)\sin\beta
= \sin\alpha\cos\beta + \cos\alpha\sin\beta$.

Тангенс: $\tan(\alpha+\beta) = \dfrac{\sin(\alpha+\beta)}{\cos(\alpha+\beta)}
= \dfrac{\sin\alpha\cos\beta + \cos\alpha\sin\beta}{\cos\alpha\cos\beta - \sin\alpha\sin\beta}$.
Делим числитель и знаменатель на $\cos\alpha\cos\beta$:
$= \dfrac{\tan\alpha + \tan\beta}{1 - \tan\alpha\tan\beta}$.
\end{proof}

%-----------------------------------------
\section{Формулы двойного угла}
%-----------------------------------------

\begin{theorem}[Формулы двойного угла]
\[
\sin 2\alpha = 2\sin\alpha\cos\alpha,
\]
\[
\cos 2\alpha = \cos^2\alpha - \sin^2\alpha = 1 - 2\sin^2\alpha
= 2\cos^2\alpha - 1,
\]
\[
\tan 2\alpha = \frac{2\tan\alpha}{1 - \tan^2\alpha}.
\]
\end{theorem}

\begin{proof}
Подставим $\beta = \alpha$ в формулы сложения:
$\sin(\alpha+\alpha) = \sin\alpha\cos\alpha + \cos\alpha\sin\alpha = 2\sin\alpha\cos\alpha$.
$\cos 2\alpha = \cos^2\alpha - \sin^2\alpha$.
Используя ОТТ: $= 1 - \sin^2\alpha - \sin^2\alpha = 1 - 2\sin^2\alpha$;
$= \cos^2\alpha - (1 - \cos^2\alpha) = 2\cos^2\alpha - 1$.
\end{proof}

%-----------------------------------------
\section{Теорема синусов}
%-----------------------------------------

\begin{theorem}[Теорема синусов]
В треугольнике $ABC$ с углами $A,B,C$ и противолежащими сторонами
$a,b,c$, вписанным в окружность радиуса $R$:
\[
\frac{a}{\sin A} = \frac{b}{\sin B} = \frac{c}{\sin C} = 2R.
\]
\end{theorem}

\begin{proof}
Докажем $a / \sin A = 2R$.

Пусть $O$ --- центр описанной окружности радиуса $R$.
Рассмотрим хорду $BC = a$. Центральный угол $\angle BOC = 2A$
(вписанный угол $A$ вдвое меньше центрального, опирающегося
на ту же дугу $BC$, не содержащую $A$).

В равнобедренном треугольнике $OBC$: $OB = OC = R$.
По теореме синусов для него:
$\dfrac{BC}{\sin(\angle BOC)} = 2R$, т.\,е.\
$\dfrac{a}{\sin 2A / ???}$...

Точнее: используем хорду. $BC = 2R\sin(\angle BOC / 2) = 2R\sin A$.
Значит $\dfrac{a}{\sin A} = 2R$.
\end{proof}

%-----------------------------------------
\section{Теорема косинусов}
%-----------------------------------------

\begin{theorem}[Теорема косинусов]
В треугольнике $ABC$:
\[
c^2 = a^2 + b^2 - 2ab\cos C.
\]
\end{theorem}

\begin{proof}
Введём координаты: $C = (0,0)$, $B = (a, 0)$,
$A = (b\cos C, b\sin C)$.
\[
c^2 = |AB|^2 = (b\cos C - a)^2 + (b\sin C)^2
= b^2\cos^2 C - 2ab\cos C + a^2 + b^2\sin^2 C
= a^2 + b^2 - 2ab\cos C.
\]
\end{proof}

%-----------------------------------------
\section{Формулы произведения}
%-----------------------------------------

\begin{theorem}[Формулы произведения]
\[
\sin\alpha\cos\beta = \tfrac{1}{2}[\sin(\alpha+\beta) + \sin(\alpha-\beta)],
\]
\[
\cos\alpha\cos\beta = \tfrac{1}{2}[\cos(\alpha-\beta) + \cos(\alpha+\beta)],
\]
\[
\sin\alpha\sin\beta = \tfrac{1}{2}[\cos(\alpha-\beta) - \cos(\alpha+\beta)].
\]
\end{theorem}

\begin{proof}
Сложим/вычтем формулы синуса суммы и разности:
$\sin(\alpha+\beta) + \sin(\alpha-\beta)
= 2\sin\alpha\cos\beta$. Делим на 2.
Аналогично для остальных.
\end{proof}

%-----------------------------------------
\section{Сумма и разность синусов и косинусов}
%-----------------------------------------

\begin{theorem}[Формулы суммы и разности]
\[
\sin\alpha + \sin\beta = 2\sin\frac{\alpha+\beta}{2}\cos\frac{\alpha-\beta}{2},
\]
\[
\sin\alpha - \sin\beta = 2\cos\frac{\alpha+\beta}{2}\sin\frac{\alpha-\beta}{2},
\]
\[
\cos\alpha + \cos\beta = 2\cos\frac{\alpha+\beta}{2}\cos\frac{\alpha-\beta}{2},
\]
\[
\cos\alpha - \cos\beta = -2\sin\frac{\alpha+\beta}{2}\sin\frac{\alpha-\beta}{2}.
\]
\end{theorem}

\begin{proof}
Обозначим $p = \dfrac{\alpha+\beta}{2}$, $q = \dfrac{\alpha-\beta}{2}$,
тогда $\alpha = p+q$, $\beta = p-q$.
\[
\sin\alpha + \sin\beta = \sin(p+q) + \sin(p-q)
= (\sin p\cos q + \cos p\sin q) + (\sin p\cos q - \cos p\sin q)
= 2\sin p\cos q.
\]
\end{proof}

%%%%%%%%%%%%%%%%%%%%%%%%%%%%%%%%%%%%%%%%%%%%%%%%%%%%%%%%%%%%%
\part{Комплексные числа}
%%%%%%%%%%%%%%%%%%%%%%%%%%%%%%%%%%%%%%%%%%%%%%%%%%%%%%%%%%%%%

\section{Бинарная операция. Свойства. Аксиомы поля.}\label{sec:field-axioms}[

\begin{definition}
\textbf{Бинарная операция} на множестве $S$ ---
отображение $*\colon S \times S \to S$.
\end{definition}

\begin{definition}
\textbf{Поле} --- множество $F$ с двумя бинарными операциями
$+$ (сложение) и $\cdot$ (умножение), удовлетворяющими следующим
аксиомам:
\begin{enumerate}[(Ф1)]
\item $(F, +)$ --- абелева группа:
      \begin{itemize}
      \item ассоциативность: $(a+b)+c = a+(b+c)$,
      \item коммутативность: $a+b = b+a$,
      \item нейтральный элемент: $\exists\, 0\colon a+0 = a$,
      \item обратный элемент: $\exists\, {-a}\colon a+(-a) = 0$.
      \end{itemize}
\item $(F\setminus\{0\}, \cdot)$ --- абелева группа:
      \begin{itemize}
      \item ассоциативность: $(a\cdot b)\cdot c = a\cdot(b\cdot c)$,
      \item коммутативность: $a\cdot b = b\cdot a$,
      \item нейтральный элемент: $\exists\, 1\neq 0\colon a\cdot 1 = a$,
      \item обратный элемент: $\forall a\neq 0\;\exists\, a^{-1}\colon a\cdot a^{-1}=1$.
      \end{itemize}
\item Дистрибутивность: $a\cdot(b+c) = a\cdot b + a\cdot c$.
\end{enumerate}
\end{definition}

%-----------------------------------------
\section{Комплексные числа. Определения. Алгебраическая форма.}
%-----------------------------------------

\begin{definition}
\textbf{Комплексное число} --- упорядоченная пара $(a,b) \in \mathbb{R}^2$.
Обозначение: $z = a + bi$, где $i$ --- \emph{мнимая единица}.

\textbf{Вещественная часть:} $\mathrm{Re}\,z = a$.
\textbf{Мнимая часть:} $\mathrm{Im}\,z = b$.
\end{definition}

\begin{definition}
\textbf{Сложение:} $(a+bi) + (c+di) = (a+c) + (b+d)i$.

\textbf{Умножение:} $(a+bi)(c+di) = (ac-bd) + (ad+bc)i$

(используя $i^2 = -1$).
\end{definition}

\begin{theorem}
Множество $\mathbb{C}$ с операциями сложения и умножения образует поле.
\end{theorem}

\begin{proof}
Проверяем аксиомы поля:
\begin{enumerate}
\item \textbf{Ассоциативность сложения:} следует из ассоциативности в $\mathbb{R}$.
\item \textbf{Коммутативность сложения:} $a+bi + c+di = (a+c)+(b+d)i =
      c+di + a+bi$.
\item \textbf{Нейтральный элемент:} $0 = 0 + 0i$.
\item \textbf{Противоположный:} $-(a+bi) = -a + (-b)i$.
\item \textbf{Коммутативность умножения:}
      $(a+bi)(c+di) = (ac-bd)+(ad+bc)i =
      (ca-db)+(da+cb)i = (c+di)(a+bi)$.
\item \textbf{Ассоциативность умножения:} прямая проверка.
\item \textbf{Нейтральный элемент:} $1 = 1 + 0i$.
\item \textbf{Обратный элемент:} для $z = a+bi \neq 0$:
      $z^{-1} = \dfrac{a - bi}{a^2+b^2}$.
      Проверка: $z \cdot z^{-1} = (a+bi)\dfrac{a-bi}{a^2+b^2}
      = \dfrac{a^2+b^2}{a^2+b^2} = 1$.
\item \textbf{Дистрибутивность:} прямая проверка.
\end{enumerate}
\end{proof}

%-----------------------------------------
\section{Сопряжение. Аргумент и модуль.}
%-----------------------------------------

\begin{definition}
\textbf{Комплексно сопряжённое} к $z = a+bi$:
$\overline{z} = a - bi$.

Геометрически: отражение относительно вещественной оси.
\end{definition}

\begin{property}[Свойства сопряжения]
\begin{enumerate}[(i)]
\item $\overline{z+w} = \overline{z} + \overline{w}$.
\item $\overline{z \cdot w} = \overline{z} \cdot \overline{w}$.
\item $z + \overline{z} = 2\,\mathrm{Re}\,z$.
\item $z \cdot \overline{z} = |z|^2$.
\item $\overline{\overline{z}} = z$.
\item $z \in \mathbb{R} \iff z = \overline{z}$.
\end{enumerate}
\end{property}

\begin{definition}
\textbf{Модуль} комплексного числа $z = a+bi$:
\[
|z| = \sqrt{a^2 + b^2}.
\]
Геометрически: расстояние от $z$ до начала координат.
\end{definition}

\begin{definition}
\textbf{Аргумент} комплексного числа $z \neq 0$ ---
угол $\varphi$ такой, что $z = |z|(\cos\varphi + i\sin\varphi)$.
Обозначение: $\arg z = \varphi$ (определён с точностью до $2\pi k$).
\textbf{Главное значение аргумента:} $\mathrm{Arg}\,z \in (-\pi, \pi]$.
\end{definition}

%-----------------------------------------
\section{Тригонометрическая форма. Умножение. Формула Муавра. Корень.}
%-----------------------------------------

\begin{definition}
\textbf{Тригонометрическая форма:} $z = r(\cos\varphi + i\sin\varphi)$,
где $r = |z|$, $\varphi = \arg z$.
\end{definition}

\begin{theorem}[Умножение в тригонометрической форме]
\[
z_1 z_2 = r_1 r_2 (\cos(\varphi_1 + \varphi_2) + i\sin(\varphi_1+\varphi_2)).
\]
\end{theorem}

\begin{proof}
$z_1 z_2 = r_1(\cos\varphi_1 + i\sin\varphi_1) \cdot r_2(\cos\varphi_2 + i\sin\varphi_2)$
$= r_1 r_2[(\cos\varphi_1\cos\varphi_2 - \sin\varphi_1\sin\varphi_2)
+ i(\sin\varphi_1\cos\varphi_2 + \cos\varphi_1\sin\varphi_2)]$
$= r_1 r_2[\cos(\varphi_1+\varphi_2) + i\sin(\varphi_1+\varphi_2)]$.
\end{proof}

\begin{theorem}[Формула Муавра]
\[
(\cos\varphi + i\sin\varphi)^n = \cos(n\varphi) + i\sin(n\varphi),
\quad n \in \mathbb{Z}.
\]
\end{theorem}

\begin{proof}
Индукция по $n \geq 0$: применяем формулу умножения.
Для $n < 0$: $z^n = (z^{-1})^{|n|}$, $z^{-1} = \cos(-\varphi)+i\sin(-\varphi)$.
\end{proof}

\begin{theorem}[Корень из комплексного числа]
Уравнение $w^n = z$ ($z \neq 0$, $n \in \mathbb{N}$) имеет ровно $n$
различных решений:
\[
w_k = \sqrt[n]{r}\Bigl(\cos\frac{\varphi + 2\pi k}{n} + i\sin\frac{\varphi+2\pi k}{n}\Bigr),
\quad k = 0, 1, \ldots, n-1.
\]
\end{theorem}

\begin{proof}
Пусть $w = \rho(\cos\psi + i\sin\psi)$. Из $w^n = z$:
$\rho^n = r$ (вещественные модули) $\Rightarrow \rho = \sqrt[n]{r}$.
$n\psi = \varphi + 2\pi k$ $\Rightarrow \psi_k = (\varphi + 2\pi k)/n$.
При $k = 0, 1, \ldots, n-1$ получаем $n$ различных значений $\psi_k \mod 2\pi$.
\end{proof}

%-----------------------------------------
\section{Корни из 1. Первообразный корень.}
%-----------------------------------------

\begin{definition}
\textbf{Корень $n$-й степени из 1} --- решение $w^n = 1$:
\[
\varepsilon_k = \cos\frac{2\pi k}{n} + i\sin\frac{2\pi k}{n},
\quad k = 0,1,\ldots,n-1.
\]
\end{definition}

\begin{property}[Свойства корней из 1]
\begin{enumerate}[(i)]
\item Корни образуют правильный $n$-угольник на единичной окружности.
\item Произведение двух корней --- корень: $\varepsilon_k \cdot \varepsilon_m = \varepsilon_{k+m \bmod n}$.
\item Они образуют \textbf{группу} по умножению, изоморфную $\mathbb{Z}/n\mathbb{Z}$.
\item Сумма всех корней: $\sum_{k=0}^{n-1} \varepsilon_k = 0$ (при $n \geq 2$).
\item $\varepsilon_k = \varepsilon_1^k$, где $\varepsilon_1 = e^{2\pi i/n}$.
\end{enumerate}
\end{property}

\begin{definition}
$\varepsilon$ --- \textbf{первообразный корень} $n$-й степени из 1,
если $\varepsilon^n = 1$ и $\varepsilon^k \neq 1$ для $0 < k < n$.
Эквивалентно: $\varepsilon = \varepsilon_k$ при $\gcd(k,n) = 1$.
\end{definition}

\begin{property}
Степени первообразного корня $\varepsilon$ перечисляют все корни $n$-й
степени из 1: $\{1, \varepsilon, \varepsilon^2, \ldots, \varepsilon^{n-1}\}$.
\end{property}

%-----------------------------------------
\section{Целые гауссовы числа}
%-----------------------------------------

\begin{definition}
\textbf{Целые гауссовы числа:}
\[
\mathbb{Z}[i] = \{a + bi \mid a, b \in \mathbb{Z}\}.
\]
Они образуют \textbf{кольцо} (замкнуты относительно $+$ и $\cdot$).
\end{definition}

\begin{example}
$3 + 2i,\; -1 + 4i,\; 5 \in \mathbb{Z}[i]$.
$(2+i)(3-i) = 6 - 2i + 3i - i^2 = 7 + i \in \mathbb{Z}[i]$.
\end{example}

%-----------------------------------------
\section{Обратимые элементы. Ассоциированность. Делимость.}
%-----------------------------------------

\begin{definition}
$u \in \mathbb{Z}[i]$ --- \textbf{обратимый} (единица), если
$\exists\, v \in \mathbb{Z}[i]\colon uv = 1$.

\textbf{Единицы $\mathbb{Z}[i]$:} $\pm 1, \pm i$.
\end{definition}

\begin{proof}
Если $uv = 1$, то $N(u)N(v) = N(1) = 1$.
Так как $N(u), N(v) \in \mathbb{Z}_{\geq 0}$, то $N(u) = N(v) = 1$.
$N(a+bi) = a^2+b^2 = 1 \iff (a,b) \in \{(\pm 1, 0),(0,\pm 1)\}$,
т.\,е.\ $u \in \{1,-1,i,-i\}$.
\end{proof}

\begin{definition}
$\alpha$ и $\beta$ в $\mathbb{Z}[i]$ \textbf{ассоциированы} ($\alpha \sim \beta$),
если $\alpha = u\beta$ для некоторой единицы $u$.
\end{definition}

\begin{definition}
$\alpha$ \textbf{делит} $\beta$ ($\alpha \mid \beta$), если
$\exists\, \gamma \in \mathbb{Z}[i]\colon \beta = \alpha\gamma$.
\end{definition}

\begin{definition}
$\alpha$ и $\beta$ \textbf{взаимно просты}, если их общие делители ---
только единицы: $\text{НОД}(\alpha,\beta) \sim 1$.
\end{definition}

%-----------------------------------------
\section{Норма гауссового числа. Свойства.}
%-----------------------------------------

\begin{definition}
\textbf{Норма} гауссового числа $z = a+bi$:
\[
N(z) = z\overline{z} = a^2 + b^2 \in \mathbb{Z}_{\geq 0}.
\]
\end{definition}

\begin{property}[Свойства нормы]
\begin{enumerate}[(i)]
\item $N(zw) = N(z)N(w)$ (мультипликативность).
\item $N(z) = 0 \iff z = 0$.
\item $N(z) = 1 \iff z$ --- единица $\mathbb{Z}[i]$.
\item $z \mid w \Rightarrow N(z) \mid N(w)$ в $\mathbb{Z}$.
\end{enumerate}
\end{property}

\begin{proof}
(i): $N(zw) = (zw)\overline{(zw)} = z\overline{z}w\overline{w} = N(z)N(w)$.
\end{proof}

%-----------------------------------------
\section{Простые гауссовы числа}
%-----------------------------------------

\begin{definition}
$\pi \in \mathbb{Z}[i]$ --- \textbf{простое гауссово число}, если:
\begin{enumerate}[(1)]
\item $\pi$ не является нулём и не единицей,
\item из $\pi = \alpha\beta$ следует, что $\alpha$ или $\beta$ --- единица.
\end{enumerate}
\end{definition}

\begin{example}
$1+i$ --- простое: $N(1+i)=2$. Если $(1+i)=\alpha\beta$, то
$N(\alpha)N(\beta)=2$, значит одно из $N(\alpha),N(\beta)$ равно 1 ---
единица. $3$ --- простое в $\mathbb{Z}[i]$. $5 = (2+i)(2-i)$ ---
не простое.
\end{example}

%-----------------------------------------
\section{Рождественская теорема Ферма}
%-----------------------------------------

\begin{theorem}[Рождественская теорема Ферма]
Простое натуральное число $p$ представляется как сумма двух квадратов
$p = a^2 + b^2$ ($a, b \in \mathbb{Z}$) тогда и только тогда, когда
$p = 2$ или $p \equiv 1 \pmod{4}$.
\end{theorem}

\begin{proof}[Идея доказательства]
\textbf{Необходимость.} Квадраты по модулю 4 равны 0 или 1.
Если $p = a^2 + b^2$:
$p \not\equiv 3 \pmod{4}$ (т.\,к.\ $0+0=0,0+1=1,1+1=2\neq 3$).
Если $p$ нечётно, то $p \equiv 1 \pmod 4$.

\textbf{Достаточность для $p \equiv 1\pmod 4$:}
\begin{enumerate}
\item По теореме Вильсона и критерию: $-1$ --- квадратичный вычет
      mod $p$ при $p \equiv 1 \pmod 4$. Существует $m\colon m^2 \equiv -1 \pmod p$.
\item $p \mid m^2 + 1 = (m+i)(m-i)$ в $\mathbb{Z}[i]$.
\item Если бы $p$ было простым в $\mathbb{Z}[i]$, то $p \mid m+i$
      или $p \mid m-i$ --- невозможно (т.\,к.\ $\mathrm{Im}(m\pm i)/p = \pm 1/p \notin \mathbb{Z}$).
\item Значит, $p$ не простое в $\mathbb{Z}[i]$: $p = \alpha\beta$,
      $N(\alpha)N(\beta) = N(p) = p^2$, $N(\alpha)=N(\beta)=p$.
      Если $\alpha = a+bi$, то $p = a^2+b^2$.
\end{enumerate}
\end{proof}

%-----------------------------------------
\section{Критерий Гаусса}
%-----------------------------------------

\begin{theorem}[Критерий Гаусса --- классификация простых гауссовых чисел]
Простые гауссовы числа (с точностью до ассоциированности) --- это:
\begin{enumerate}[(1)]
\item $1 + i$ (единственное над $p=2$, т.\,к.\ $2 = -i(1+i)^2$),
\item $a + bi$ с $a^2 + b^2 = p$, где $p \equiv 1 \pmod 4$ --- простое в $\mathbb{Z}$,
\item $q$ --- нечётное простое в $\mathbb{Z}$ с $q \equiv 3 \pmod 4$.
\end{enumerate}
\end{theorem}

\begin{proof}[Идея]
Пусть $\pi \in \mathbb{Z}[i]$ --- простое гауссово число.
Тогда $N(\pi) = p$ или $N(\pi) = p^2$ для некоторого простого $p \in \mathbb{Z}$.

Случай $N(\pi) = p$: $\pi\overline{\pi} = p$, значит $p$ раскладывается
в $\mathbb{Z}[i]$. По РТФ, $p = 2$ или $p \equiv 1 \pmod 4$.

Случай $N(\pi) = p^2$: $\pi = p$ (простое в $\mathbb{Z}$, $p\equiv 3\pmod 4$).
\end{proof}

%-----------------------------------------
\section{ОТА (Основная теорема арифметики) в $\mathbb{Z}[i]$}
%-----------------------------------------

\begin{theorem}[ОТА в $\mathbb{Z}\lbrack i\rbrack$]
Кольцо $\mathbb{Z}[i]$ --- UFD (кольцо факториального разложения):
каждый ненулевой ненеобратимый элемент представим в виде произведения
простых, и это представление единственно с точностью до порядка
и ассоциированности.
\end{theorem}

\begin{proof}[Идея]
$\mathbb{Z}[i]$ --- евклидово кольцо с нормой $N$. Евклидово кольцо
является областью главных идеалов (ОГИ), а ОГИ --- UFD.
Евклидовость: для $\alpha, \beta \neq 0$ существуют $\gamma, \rho$ такие,
что $\alpha = \beta\gamma + \rho$ и $N(\rho) < N(\beta)$
(деление с остатком).
\end{proof}

%-----------------------------------------
\section{Деление с остатком в $\mathbb{Z}[i]$}
%-----------------------------------------

\begin{theorem}[Деление с остатком]
Для любых $\alpha, \beta \in \mathbb{Z}[i]$, $\beta \neq 0$,
существуют $\gamma, \rho \in \mathbb{Z}[i]$ такие, что
\[
\alpha = \beta\gamma + \rho, \quad N(\rho) < N(\beta).
\]
\end{theorem}

\begin{proof}
Рассмотрим $\alpha/\beta = x + yi \in \mathbb{C}$.
Выберем $\gamma = m + ni$, где $m = \mathrm{round}(x)$,
$n = \mathrm{round}(y)$ (ближайшие целые).
Тогда $|x - m| \leq 1/2$, $|y - n| \leq 1/2$.
$\rho = \alpha - \beta\gamma$.
$N(\rho/\beta) = N(\alpha/\beta - \gamma) = (x-m)^2 + (y-n)^2 \leq 1/4 + 1/4 = 1/2 < 1$.
$N(\rho) = N(\rho/\beta) \cdot N(\beta) < N(\beta)$.
\end{proof}

%-----------------------------------------
\section{НОД. Алгоритм Евклида. Линейное представление.}
%-----------------------------------------

\begin{definition}
$\delta \in \mathbb{Z}[i]$ --- \textbf{НОД} $\alpha$ и $\beta$, если:
\begin{enumerate}[(1)]
\item $\delta \mid \alpha$ и $\delta \mid \beta$,
\item если $\gamma \mid \alpha$ и $\gamma \mid \beta$, то $\gamma \mid \delta$.
\end{enumerate}
НОД определён с точностью до ассоциированности.
\end{definition}

\begin{method}[Алгоритм Евклида для $\mathbb{Z}\lbrack i\rbrack$]
\begin{align*}
\alpha &= \beta q_1 + r_1, \quad N(r_1) < N(\beta), \\
\beta &= r_1 q_2 + r_2, \quad N(r_2) < N(r_1), \\
&\;\vdots\\
r_{k-1} &= r_k q_{k+1} + 0.
\end{align*}
Тогда $\gcd(\alpha,\beta) \sim r_k$.
\end{method}

\begin{theorem}[Линейное представление НОД]
Существуют $s, t \in \mathbb{Z}[i]$ такие, что
$\gcd(\alpha,\beta) = s\alpha + t\beta$.
\end{theorem}

\begin{proof}
Обратный ход алгоритма Евклида (расширенный алгоритм Евклида).
\end{proof}

%-----------------------------------------
\section{Геометрия гауссовых чисел}
%-----------------------------------------

\begin{note}
Числа, кратные $\alpha = a+bi$, образуют подрешётку $\mathbb{Z}[i]$:
$\{\alpha k \mid k \in \mathbb{Z}[i]\}$.
Фундаментальный параллелограмм имеет площадь $N(\alpha) = a^2+b^2$.
\end{note}

\begin{note}
Числа, равные $|\alpha|$ по модулю, лежат на окружности
$x^2 + y^2 = N(\alpha)$.
Гауссовы числа на этой окружности ---
решения $a'^2 + b'^2 = N(\alpha)$.
\end{note}

%-----------------------------------------
\section{Уравнение прямой на комплексной плоскости}
%-----------------------------------------

\begin{theorem}[Уравнение прямой]
Прямая через $z_1, z_2 \in \mathbb{C}$ задаётся условием:
\[
\mathrm{Im}\!\left(\frac{z - z_1}{\overline{z_2 - z_1}}\right) = 0,
\]
что эквивалентно:
\[
\begin{vmatrix} z & \overline{z} & 1 \\ z_1 & \overline{z_1} & 1 \\
z_2 & \overline{z_2} & 1 \end{vmatrix} = 0.
\]
\end{theorem}

\begin{definition}
Четыре точки $z_1,z_2,z_3,z_4$ лежат на одной окружности или прямой
тогда и только тогда, когда их \textbf{двойное отношение}
\[
(z_1,z_2;z_3,z_4) = \frac{(z_1-z_3)(z_2-z_4)}{(z_1-z_4)(z_2-z_3)}
\]
является вещественным.
\end{definition}

%-----------------------------------------
\section{Параллельные и перпендикулярные прямые}
%-----------------------------------------

\begin{theorem}
Прямые через $z_1,z_2$ и через $w_1,w_2$:
\begin{itemize}
\item \textbf{параллельны} $\iff$ $\dfrac{z_2-z_1}{w_2-w_1} \in \mathbb{R}$
      (аргументы разности одинаковы),
\item \textbf{перпендикулярны} $\iff$ $\dfrac{z_2-z_1}{w_2-w_1} \in i\mathbb{R}$
      (аргументы разности отличаются на $\pi/2$).
\end{itemize}
\end{theorem}

%-----------------------------------------
\section{Углы в комплексной геометрии. Вписанный четырёхугольник.}
%-----------------------------------------

\begin{definition}
\textbf{Угол} от луча $[z_0,z_1)$ до луча $[z_0,z_2)$ ---
аргумент дроби $\dfrac{z_2 - z_0}{z_1 - z_0}$.
\end{definition}

\begin{theorem}[Вписанный четырёхугольник]
Четыре точки $A,B,C,D$ лежат на одной окружности тогда и только тогда,
когда
\[
\arg\frac{C - A}{B - A} = \arg\frac{C - D}{B - D} \pmod{\pi},
\]
т.\,е.\ вписанные углы $\angle CAB = \angle CDB$.
Эквивалентно: двойное отношение $(A,B;C,D) \in \mathbb{R}$.
\end{theorem}

%%%%%%%%%%%%%%%%%%%%%%%%%%%%%%%%%%%%%%%%%%%%%%%%%%%%%%%%%%%%%
\part{Линейная алгебра}
%%%%%%%%%%%%%%%%%%%%%%%%%%%%%%%%%%%%%%%%%%%%%%%%%%%%%%%%%%%%%

\section{Аксиомы поля}

(См.~часть IV, \S\ref{sec:field-axioms} --- полностью изложено там.)

%-----------------------------------------
\section{Векторное пространство и подпространство}
%-----------------------------------------

\begin{definition}
\textbf{Векторное пространство} над полем $F$ --- множество $V$
с операциями:
\begin{itemize}
\item сложение: $+ \colon V \times V \to V$,
\item умножение на скаляр: $\cdot \colon F \times V \to V$,
\end{itemize}
удовлетворяющими аксиомам:
\begin{enumerate}[(В1)]
\item $(V, +)$ --- абелева группа (аксиомы как в поле для $+$).
\item $1 \cdot v = v$.
\item $(\lambda\mu)v = \lambda(\mu v)$.
\item $(\lambda + \mu)v = \lambda v + \mu v$.
\item $\lambda(u + v) = \lambda u + \lambda v$.
\end{enumerate}
\end{definition}

\begin{example}
$\mathbb{R}^n$, $F^n$, пространство многочленов $F[x]$,
пространство матриц $M_{m\times n}(F)$, пространство функций.
\end{example}

\begin{definition}
Непустое подмножество $W \subseteq V$ --- \textbf{подпространство},
если оно замкнуто относительно сложения и умножения на скаляр:
\[
u, v \in W,\; \lambda \in F \Rightarrow u + v \in W,\; \lambda v \in W.
\]
\end{definition}

\begin{example}
$\{0\}$, $V$, множество решений однородной СЛАУ, пространство чётных функций.
\end{example}

%-----------------------------------------
\section{Линейная зависимость. Элементарные преобразования.}
%-----------------------------------------

\begin{definition}
Векторы $v_1, \ldots, v_k \in V$ \textbf{линейно зависимы}, если
существуют $\lambda_1,\ldots,\lambda_k \in F$ (не все нулевые) такие, что
$\lambda_1 v_1 + \cdots + \lambda_k v_k = 0$.
В противном случае --- \textbf{линейно независимы}.
\end{definition}

\begin{definition}
\textbf{Элементарные преобразования} строк матрицы:
\begin{enumerate}[(Э1)]
\item Перестановка двух строк.
\item Умножение строки на ненулевой скаляр.
\item Прибавление к одной строке другой, умноженной на скаляр.
\end{enumerate}
Элементарные преобразования не меняют ранг матрицы и множество
решений СЛАУ (при применении к расширенной матрице).
\end{definition}

%-----------------------------------------
\section{Линейная оболочка, порождающая система, базис}
%-----------------------------------------

\begin{definition}
\textbf{Линейная оболочка} (span) множества $S = \{v_1,\ldots,v_k\}$:
\[
\mathrm{span}(S) = \Bigl\{\sum_{i=1}^k \lambda_i v_i \Bigm| \lambda_i \in F\Bigr\}.
\]
Это наименьшее подпространство, содержащее $S$.
\end{definition}

\begin{definition}
$S$ --- \textbf{порождающая система} $V$, если $\mathrm{span}(S) = V$.
\end{definition}

\begin{definition}
\textbf{Базис} пространства $V$ --- линейно независимая порождающая система.
\end{definition}

%-----------------------------------------
\section{Теоремы про базис}
%-----------------------------------------

\begin{theorem}[Существование базиса]
Каждое конечномерное векторное пространство $V \neq \{0\}$ имеет базис.
\end{theorem}

\begin{proof}
Из любой порождающей системы можно удалять векторы, выражающиеся
через оставшиеся, пока не получим линейно независимую систему.
Она остаётся порождающей.
\end{proof}

\begin{theorem}[Единственность разложения по базису]
Если $e_1,\ldots,e_n$ --- базис $V$, то каждый вектор $v \in V$
единственным образом представляется в виде
$v = \lambda_1 e_1 + \cdots + \lambda_n e_n$.
\end{theorem}

\begin{proof}
\textbf{Существование:} по определению порождающей системы.
\textbf{Единственность:} если $v = \sum \lambda_i e_i = \sum \mu_i e_i$,
то $\sum (\lambda_i - \mu_i) e_i = 0$, откуда $\lambda_i = \mu_i$ по
линейной независимости.
\end{proof}

\begin{theorem}[Фиксированный размер базиса]
Все базисы конечномерного векторного пространства $V$ имеют
одинаковое количество элементов --- \textbf{размерность} $\dim V$.
\end{theorem}

\begin{proof}
Пусть $e_1,\ldots,e_n$ и $f_1,\ldots,f_m$ --- два базиса.
Каждый $f_j$ выражается через $e_i$. Матрица перехода $n\times m$.
Если $m > n$, то система $m$ уравнений в $n$ переменных имеет
нетривиальное решение --- $f_j$ линейно зависимы. Противоречие.
Аналогично $n \leq m$. Значит $n = m$.
\end{proof}

%-----------------------------------------
\section{Матрицы. Основные понятия, арифметические операции.}
%-----------------------------------------

\begin{definition}
\textbf{Матрица} размера $m \times n$ над полем $F$ ---
прямоугольная таблица элементов $A = (a_{ij})$,
$i = 1,\ldots,m$, $j = 1,\ldots,n$.
\end{definition}

\begin{definition}
\textbf{Сложение} ($A + B$ при одинаковых размерах):
$(A+B)_{ij} = a_{ij} + b_{ij}$.

\textbf{Умножение на скаляр:} $(\lambda A)_{ij} = \lambda a_{ij}$.

\textbf{Умножение матриц} ($A\colon m\times n$, $B\colon n\times p$):
$(AB)_{ij} = \sum_{k=1}^n a_{ik} b_{kj}$.
\end{definition}

\begin{property}
\begin{enumerate}[(i)]
\item Умножение ассоциативно: $(AB)C = A(BC)$.
\item Дистрибутивность: $A(B+C) = AB + AC$.
\item Умножение \emph{не} коммутативно в общем случае: $AB \neq BA$.
\end{enumerate}
\end{property}

%-----------------------------------------
\section{Транспонирование матриц}
%-----------------------------------------

\begin{definition}
\textbf{Транспонированная матрица} $A^T$ размера $n \times m$:
$(A^T)_{ij} = a_{ji}$.
\end{definition}

\begin{property}
\begin{enumerate}[(i)]
\item $(A^T)^T = A$.
\item $(A + B)^T = A^T + B^T$.
\item $(\lambda A)^T = \lambda A^T$.
\item $(AB)^T = B^T A^T$.
\end{enumerate}
\end{property}

\begin{proof}[Доказательство (iv)]
$((AB)^T)_{ij} = (AB)_{ji} = \sum_k a_{jk}b_{ki}
= \sum_k (B^T)_{ik}(A^T)_{kj} = (B^T A^T)_{ij}$.
\end{proof}

%-----------------------------------------
\section{Перестановки}
%-----------------------------------------

\begin{definition}
\textbf{Перестановка} множества $\{1,\ldots,n\}$ --- биекция
$\sigma\colon \{1,\ldots,n\} \to \{1,\ldots,n\}$.
Множество всех перестановок: $S_n$ (симметрическая группа), $|S_n| = n!$.
\end{definition}

\begin{definition}
\textbf{Транспозиция} --- перестановка, меняющая местами два элемента:
$(ij)(\sigma) = \sigma$ с $i$ и $j$ поменянными местами.
\end{definition}

\begin{definition}
\textbf{Цикл} $(a_1\, a_2\,\cdots\, a_k)$ --- перестановка:
$a_1 \mapsto a_2 \mapsto \cdots \mapsto a_k \mapsto a_1$,
остальные на месте.
\end{definition}

\begin{definition}
\textbf{Инверсия} перестановки $\sigma$ --- пара $(i,j)$ с $i < j$,
$\sigma(i) > \sigma(j)$.
Число инверсий: $\mathrm{inv}(\sigma)$.
\end{definition}

\begin{definition}
\textbf{Знак перестановки:}
$\mathrm{sgn}(\sigma) = (-1)^{\mathrm{inv}(\sigma)}$.
$\sigma$ --- \textbf{чётная}, если $\mathrm{sgn}(\sigma) = 1$;
\textbf{нечётная}, если $\mathrm{sgn}(\sigma) = -1$.
\end{definition}

\begin{property}
$\mathrm{sgn}(\sigma\tau) = \mathrm{sgn}(\sigma)\cdot\mathrm{sgn}(\tau)$.
Транспозиция --- нечётная перестановка.
\end{property}

%-----------------------------------------
\section{Определитель. Свойства.}
%-----------------------------------------

\begin{definition}
\textbf{Определитель} (детерминант) матрицы $A = (a_{ij})$ размера $n\times n$:
\[
\det A = \sum_{\sigma \in S_n} \mathrm{sgn}(\sigma)\,
a_{1\sigma(1)} a_{2\sigma(2)} \cdots a_{n\sigma(n)}.
\]
\end{definition}

\begin{property}
\begin{enumerate}[(i)]
\item \textbf{Полилинейность по строкам/столбцам.}
\item \textbf{Кососимметричность:} перестановка двух строк меняет знак.
\item $\det(I) = 1$.
\item $\det(AB) = \det A \cdot \det B$.
\item $\det(A^T) = \det A$.
\item $\det(\lambda A) = \lambda^n \det A$.
\item $A$ обратима $\iff \det A \neq 0$.
\item Если строка (столбец) --- линейная комбинация других, то $\det = 0$.
\item Прибавление к строке другой строки, умноженной на константу,
      не меняет определитель.
\end{enumerate}
\end{property}

\begin{proof}[$\det(AB)=\det A \cdot \det B$ (идея)]
Оба отображения $B \mapsto \det(AB)$ и $B \mapsto \det A \cdot \det B$
являются полилинейными кососимметричными функциями строк $B$,
принимающими одинаковое значение на единичной матрице.
По единственности такого функционала (нормированного) они совпадают.
\end{proof}

%-----------------------------------------
\section{Обратные матрицы. Методы поиска.}
%-----------------------------------------

\begin{definition}
Матрица $B$ --- \textbf{обратная} к $A$ ($B = A^{-1}$), если
$AB = BA = I$.
$A$ обратима $\iff \det A \neq 0$.
\end{definition}

\begin{method}[Через элементарные преобразования]
Приписываем к $A$ единичную матрицу: $(A \mid I)$ и применяем
элементарные преобразования строк до вида $(I \mid A^{-1})$.
\end{method}

\begin{method}[Через матрицу алгебраических дополнений]
\[
A^{-1} = \frac{1}{\det A} \cdot \hat{A},
\]
где $\hat{A}_{ij} = (-1)^{i+j} M_{ji}$ --- матрица алгебраических
дополнений ($M_{ji}$ --- дополнительный минор).
\end{method}

%-----------------------------------------
\section{Ранг матрицы}
%-----------------------------------------

\begin{definition}
\textbf{Ранг} матрицы $A$ --- максимальное число линейно независимых
строк (= максимальное число линейно независимых столбцов =
наибольший порядок ненулевого минора). Обозначение: $\mathrm{rk}(A)$.
\end{definition}

\begin{theorem}[Эквивалентность определений ранга]
Следующие числа равны:
\begin{enumerate}[(1)]
\item Максимальное число ЛНЗ строк.
\item Максимальное число ЛНЗ столбцов.
\item Наибольший порядок ненулевого минора.
\item Число ненулевых строк в ступенчатой форме.
\end{enumerate}
\end{theorem}

\begin{proof}[Идея]
Элементарные преобразования не меняют ранг (биекция пространства строк).
Приводим к ступенчатой форме; ненулевые строки ЛНЗ. Их число равно
размерности пространства строк. Пространства строк и столбцов равноразмерны
(фундаментальная теорема ЛА).
\end{proof}

\begin{property}
\begin{enumerate}[(i)]
\item $\mathrm{rk}(AB) \leq \min(\mathrm{rk}(A),\mathrm{rk}(B))$.
\item $\mathrm{rk}(A+B) \leq \mathrm{rk}(A) + \mathrm{rk}(B)$.
\item $\mathrm{rk}(A^T) = \mathrm{rk}(A)$.
\end{enumerate}
\end{property}

%-----------------------------------------
\section{Метод Крамера}
%-----------------------------------------

\begin{theorem}[Крамер]
Система $Ax = b$ ($A$ --- квадратная, $\det A \neq 0$) имеет
единственное решение:
\[
x_j = \frac{\det A_j}{\det A},
\]
где $A_j$ --- матрица $A$, в которой $j$-й столбец заменён на $b$.
\end{theorem}

\begin{proof}
Из $Ax = b$: $x = A^{-1}b = \dfrac{1}{\det A}\hat{A}b$.
$(x)_j = \dfrac{1}{\det A}\sum_i \hat{A}_{ji} b_i
= \dfrac{\det A_j}{\det A}$
(разложение определителя $A_j$ по $j$-му столбцу).
\end{proof}

%-----------------------------------------
\section{Теорема Кронекера--Капелли}
%-----------------------------------------

\begin{theorem}[Кронекер--Капелли]
Система $Ax = b$ совместна (имеет решение) тогда и только тогда, когда
\[
\mathrm{rk}(A) = \mathrm{rk}(A \mid b).
\]
\end{theorem}

\begin{proof}
$Ax = b$ совместна $\iff$ $b$ лежит в линейной оболочке столбцов $A$
$\iff$ добавление $b$ не увеличивает ранг матрицы столбцов.
\end{proof}

%-----------------------------------------
\section{Матрица как линейное отображение}
%-----------------------------------------

\begin{definition}
\textbf{Линейное отображение} $\varphi\colon V \to W$ --- отображение,
удовлетворяющее:
\[
\varphi(u + v) = \varphi(u) + \varphi(v),\quad
\varphi(\lambda v) = \lambda \varphi(v).
\]
\end{definition}

\begin{theorem}
Каждое линейное отображение $\varphi\colon F^n \to F^m$ задаётся
матрицей $A \in M_{m\times n}(F)$: $\varphi(x) = Ax$.
При выборе базисов $e$ в $V$ и $f$ в $W$ матрица $A$ имеет
$j$-й столбец --- образ $j$-го базисного вектора:
$Ae_j = \varphi(e_j)$.
\end{theorem}

\begin{proof}
$\varphi(x) = \varphi\bigl(\sum_j x_j e_j\bigr) = \sum_j x_j \varphi(e_j)$.
Если $\varphi(e_j) = \sum_i a_{ij} f_i$, то
$\varphi(x) = \sum_i \bigl(\sum_j a_{ij} x_j\bigr) f_i = Ax$ (в координатах).
\end{proof}

%-----------------------------------------
\section{Ядро и образ линейного отображения}
%-----------------------------------------

\begin{definition}
Для линейного отображения $\varphi\colon V \to W$:
\[
\ker\varphi = \{v \in V \mid \varphi(v) = 0\} \subseteq V
\quad\text{(ядро)},
\]
\[
\mathrm{Im}\,\varphi = \{\varphi(v) \mid v \in V\} \subseteq W
\quad\text{(образ)}.
\]
\end{definition}

\begin{theorem}[Теорема о ранге и дефекте]
\[
\dim(\ker\varphi) + \dim(\mathrm{Im}\,\varphi) = \dim V.
\]
\end{theorem}

\begin{proof}
Пусть $\dim V = n$, $\dim(\ker\varphi) = k$.
Выберем базис $e_1,\ldots,e_k$ ядра и дополним его до базиса $V$:
$e_1,\ldots,e_k, e_{k+1},\ldots,e_n$.
Утверждаем, что $\varphi(e_{k+1}),\ldots,\varphi(e_n)$ --- базис образа.
\textbf{Порождение:} $\varphi(v) = \sum \lambda_i \varphi(e_i)
= \sum_{i=k+1}^n \lambda_i \varphi(e_i)$ (т.\,к.\ $\varphi(e_i)=0$
для $i \leq k$).
\textbf{ЛНЗ:} если $\sum_{i=k+1}^n \lambda_i \varphi(e_i) = 0$,
то $\varphi(\sum_{i=k+1}^n \lambda_i e_i) = 0$,
значит $\sum_{i=k+1}^n \lambda_i e_i \in \ker\varphi$,
т.\,е.\ $\sum_{i=k+1}^n \lambda_i e_i = \sum_{i=1}^k \mu_i e_i$,
откуда $\lambda_i = 0$ (т.\,к.\ базис).
Итого $\dim(\mathrm{Im}) = n-k$.
\end{proof}

%-----------------------------------------
\section{Замена базиса}
%-----------------------------------------

\begin{definition}
Пусть $e = (e_1,\ldots,e_n)$ и $e' = (e_1',\ldots,e_n')$ --- два базиса $V$.
\textbf{Матрица перехода} $P$ от $e$ к $e'$: $j$-й столбец ---
координаты $e_j'$ в базисе $e$.
\[
e_j' = \sum_i P_{ij} e_i.
\]
\end{definition}

\begin{theorem}
Если координатный вектор $v$ в базисе $e$ --- $x$, в базисе $e'$ --- $x'$,
то $x = Px'$ (т.\,е.\ $x' = P^{-1}x$).
Матрица линейного отображения $\varphi$ в базисе $e'$:
$A' = P^{-1}AP$.
\end{theorem}

%-----------------------------------------
\section{Подобные матрицы}
%-----------------------------------------

\begin{definition}
Матрицы $A$ и $B$ \textbf{подобны}, если существует обратимая $P$:
$B = P^{-1}AP$.
\end{definition}

\begin{property}
Подобные матрицы имеют одинаковые:
\begin{enumerate}[(i)]
\item определитель ($\det B = \det(P^{-1}AP) = \det A$),
\item след ($\mathrm{tr}(B) = \mathrm{tr}(A)$),
\item характеристический многочлен,
\item собственные значения.
\end{enumerate}
\end{property}

\begin{proof}
$\det(P^{-1}AP) = \det(P^{-1})\det A\det P = \frac{1}{\det P}\det A\det P = \det A$.
$\mathrm{tr}(P^{-1}AP) = \mathrm{tr}(AP P^{-1}) = \mathrm{tr}(A)$.
\end{proof}

%-----------------------------------------
\section{Теорема Лапласа. Следствия.}
%-----------------------------------------

\begin{definition}
\textbf{Алгебраическое дополнение} элемента $a_{ij}$:
$A_{ij} = (-1)^{i+j} M_{ij}$, где $M_{ij}$ --- минор
(определитель матрицы без $i$-й строки и $j$-го столбца).
\end{definition}

\begin{theorem}[Разложение Лапласа по строке/столбцу]
\[
\det A = \sum_{j=1}^n a_{ij} A_{ij}
\quad\text{(разложение по $i$-й строке)},
\]
\[
\det A = \sum_{i=1}^n a_{ij} A_{ij}
\quad\text{(разложение по $j$-му столбцу)}.
\]
\end{theorem}

\begin{proof}
Из определения определителя:
$\det A = \sum_{\sigma} \mathrm{sgn}(\sigma) \prod a_{k\sigma(k)}$.
Фиксируем $i$, группируем по $\sigma(i) = j$:
$\det A = \sum_j a_{ij} \sum_{\sigma:\sigma(i)=j} \mathrm{sgn}(\sigma)
\prod_{k\neq i} a_{k\sigma(k)}$.
Внутренняя сумма равна $(-1)^{i+j} M_{ij} = A_{ij}$.
\end{proof}

\begin{corollary}
$\sum_j a_{ij} A_{kj} = 0$ при $i \neq k$
(разложение определителя матрицы с двумя одинаковыми строками).
\end{corollary}

\begin{corollary}
$A^{-1} = \dfrac{1}{\det A}(A_{ji})$ (матрица алгебраических дополнений, транспонированная).
\end{corollary}

%-----------------------------------------
\section{Метод Гаусса}
%-----------------------------------------

\begin{method}[Метод Гаусса (Гауссово исключение)]
Для системы $Ax = b$:
\begin{enumerate}
\item Составляем расширенную матрицу $(A \mid b)$.
\item Элементарными преобразованиями строк приводим к
      \textbf{ступенчатой форме} (треугольной).
\item Обратным ходом находим решение.
\end{enumerate}
Если при приведении к ступенчатой форме в строке получается
$(0\,\cdots\,0 \mid c)$ с $c \neq 0$ --- система несовместна.
Если есть свободные переменные --- бесконечно много решений.
\end{method}

%-----------------------------------------
\section{Декартово произведение векторных пространств}
%-----------------------------------------

\begin{definition}
\textbf{Прямое произведение} (декартово произведение) пространств
$V$ и $W$ над $F$:
\[
V \times W = \{(v, w) \mid v \in V, w \in W\}
\]
с операциями:
$(v_1,w_1)+(v_2,w_2) = (v_1+v_2,\, w_1+w_2)$,
$\lambda(v,w) = (\lambda v, \lambda w)$.
\end{definition}

\begin{property}
$\dim(V \times W) = \dim V + \dim W$.
Если $(e_i)$ --- базис $V$, $(f_j)$ --- базис $W$, то
$(e_i,0)$ и $(0,f_j)$ образуют базис $V \times W$.
\end{property}

%-----------------------------------------
\section{Факторпространство}
%-----------------------------------------

\begin{definition}
Пусть $W \subseteq V$ --- подпространство.
\textbf{Смежный класс} вектора $v$: $v + W = \{v + w \mid w \in W\}$.
\textbf{Факторпространство} $V/W$ --- множество всех смежных классов
с операциями:
$(v_1+W) + (v_2+W) = (v_1+v_2)+W$,
$\lambda(v+W) = (\lambda v)+W$.
\end{definition}

\begin{theorem}
$V/W$ является векторным пространством, и $\dim(V/W) = \dim V - \dim W$.
\end{theorem}

\begin{proof}
Операции корректно определены (не зависят от выбора представителя).
Если $(e_1,\ldots,e_k)$ --- базис $W$, дополненный до базиса $V$:
$(e_1,\ldots,e_k,e_{k+1},\ldots,e_n)$, то классы
$e_{k+1}+W,\ldots,e_n+W$ образуют базис $V/W$.
\end{proof}

%-----------------------------------------
\section{Теорема об изоморфизме конечномерных пространств}
%-----------------------------------------

\begin{theorem}
Два конечномерных векторных пространства над одним полем $F$
изоморфны тогда и только тогда, когда они имеют одинаковую размерность.
\end{theorem}

\begin{proof}
\textbf{Если $\dim V = \dim W = n$:} выбираем базисы
$e_1,\ldots,e_n$ в $V$ и $f_1,\ldots,f_n$ в $W$,
определяем $\varphi(e_i) = f_i$ и продолжаем линейно.
$\varphi$ --- изоморфизм.

\textbf{Если $V \cong W$:} изоморфизм $\varphi$ переводит базис в базис,
сохраняя размерность.
\end{proof}

\begin{corollary}
Любое $n$-мерное пространство над $F$ изоморфно $F^n$.
\end{corollary}

%-----------------------------------------
\section{Теорема о гомоморфизме}
%-----------------------------------------

\begin{theorem}[Первая теорема об изоморфизме (о гомоморфизме)]
Пусть $\varphi\colon V \to W$ --- линейное отображение.
Тогда:
\[
V / \ker\varphi \;\cong\; \mathrm{Im}\,\varphi.
\]
\end{theorem}

\begin{proof}
Определим $\tilde\varphi\colon V/\ker\varphi \to \mathrm{Im}\,\varphi$ по правилу
$\tilde\varphi(v + \ker\varphi) = \varphi(v)$.

\textbf{Корректность:} если $v - v' \in \ker\varphi$, то
$\varphi(v) - \varphi(v') = \varphi(v-v') = 0$.

\textbf{Линейность:} проверяется напрямую.

\textbf{Инъективность:} $\tilde\varphi(v + \ker\varphi) = 0
\Rightarrow \varphi(v) = 0 \Rightarrow v \in \ker\varphi$,
т.\,е.\ $v + \ker\varphi = 0 + \ker\varphi$.

\textbf{Сюръективность:} $\forall w \in \mathrm{Im}\,\varphi\;
\exists v\colon \varphi(v) = w$, значит $\tilde\varphi(v+\ker\varphi) = w$.

Итого $\tilde\varphi$ --- изоморфизм.
\end{proof}

\begin{corollary}
$\dim(\mathrm{Im}\,\varphi) = \dim V - \dim(\ker\varphi)$
(теорема о ранге и дефекте).
\end{corollary}
\end{document}